\chapter{Apprentissage profond sur la variété des matrices SPD}
\label{ch:spdnet}

\section{Introduction}
\label{sec:spdnet-introduction}

\begin{contributions}
  Ce chapitre rend compte des travaux suivants, dont les versions publiées sont
  reproduites en fin de chapitre:
  \begin{itemize}
    \item \fullcite{JSTARS25};
    \item \fullcite{IGARSS24};
    \item \fullcite{PREP26b};
    \item \fullcite{EUSIPCO26b}.
  \end{itemize}
\end{contributions}

% ============================ TODO texte (~4 paragraphes) ===========================
% Introduction du chapitre. Ordre suggere:
%  1. Les trois collaborations, avec les statuts exacts --- c'est le seul endroit
%     du chapitre ou ils apparaissent, et ils different:
%       - Douba Jafuno: these CIFRE Geolithe 2022--2025, encadrement 45 %, avec
%         Guillaume Ginolhac et Nickolas Stelzenmuller (cf. le CV,
%         \cref{sec:cv-annexe-encadrements}). Soutenue le 4 novembre 2025;
%       - Matthieu Gallet: ancien stagiaire de master 2 (2020) puis doctorant,
%         aujourd'hui MCF --- la collaboration sur la batch normalization est
%         donc *posterieure* a l'encadrement, entre collegues. Ne pas la
%         presenter comme un encadrement;
%       - Thibault Pautrel: stage de master 2 en 2025 (encadrement 33 %, avec
%         Florent Bouchard et Guillaume Ginolhac), devenu post-doctorant sur le
%         meme sujet. C'est le travail le plus recent du memoire.
%  2. Le fil du chapitre, a annoncer sans detour: un reseau operant sur des
%     matrices SPD est une chaine d'applications sur une variete; a trois
%     endroits --- la donnee en amont, l'interieur du reseau, les parametres en
%     aval --- une operation standard repose sur une *moyenne euclidienne*, et
%     chacune doit ceder la place a une moyenne de Frechet
%     (\cref{def:context-riemann-frechet}).
%  3. Les deux enseignements transversaux, qui sont la these du chapitre:
%       (a) choisir la moyenne, c'est choisir une geometrie, donc un modele
%           statistique --- point developpe en
%           \cref{sec:spdnet-batchnorm-moyennes}, et qui transforme un resultat
%           experimental (« GAH bat la moyenne geometrique ») en un enonce de
%           modelisation;
%       (b) l'objet geometriquement juste est toujours iteratif, donc trop cher;
%           le travail consiste a choisir l'approximation en forme close *et a
%           en deriver le gradient a la main* --- gain d'un facteur 2 en memoire
%           en \cref{sec:spdnet-batchnorm-resultats}, et echec pur et simple de
%           la differentiation automatique sur deux des trois jeux de donnees.
%  4. Annoncer que \cref{sec:spdnet-covpool} est la seule des trois sections a
%     *mesurer* la robustesse que l'invariance geometrique promet (donnees
%     rares, etiquettes erronees, changement de conditions d'acquisition), et
%     que c'est a ce titre qu'elle ouvre le chapitre plutot qu'elle ne le clot.
% ====================================================================================

% ============================ TODO texte (~1 paragraphe) ============================
% Paragraphe de renvoi vers la partie de contexte, apres l'introduction. Meme
% principe qu'au \cref{ch:robust-detection}: le chapitre doit pouvoir se lire
% sans y revenir pour un lecteur familier. Points d'appui a mentionner:
%   - le cone et sa metrique affine invariante
%     (\cref{sec:context-riemann-spd,def:context-riemann-metriquenaturelle});
%   - la moyenne de Frechet (\cref{def:context-riemann-frechet}) et l'absence de
%     forme close pour la metrique affine invariante
%     (\cref{eq:context-riemann-frechetopt});
%   - la metrique log-euclidienne (\cref{def:context-riemann-logeuclidienne});
%   - la variete de Stiefel et la projection sur son espace tangent
%     (\cref{def:context-riemann-stiefel,eq:context-riemann-projstiefel}), et la
%     remarque sur la compacite et les retractions
%     (\cref{rem:context-riemann-stiefelvscone});
%   - le gradient riemannien et la descente
%     (\cref{sec:context-riemann-optimization});
%   - la metrique affine invariante comme geometrie de Fisher du modele gaussien
%     (\cref{sec:context-geometry-gaussian}) et les divergences
%     (\cref{sec:context-geometry-divergences}) --- ce dernier renvoi est celui
%     qui porte le plus de poids, cf. \cref{prop:spdnet-moyennes-frechet};
%   - le regime dimensionnel et le retrecissement
%     (\cref{sec:context-covariance-regularized-estimation}), qui eclairent a la
%     fois le seuil de \ReEig{} et le conditionnement des donnees.
% ====================================================================================

% ================================ TODO decision =====================================
% Ajouts a la partie contexte identifies en construisant ce chapitre, et
% volontairement non faits pour l'instant. Chacun rendrait une partie de ce
% chapitre plus courte et mieux ancree:
%  1. \cref{sec:context-geometry-divergences}: sous-section « moyennes associees
%     aux divergences » (KL a gauche -> arithmetique, KL a droite -> harmonique,
%     KL symetrisee -> GAH), avec la lecture Wishart / Wishart inverse. Sans
%     elle, \cref{prop:spdnet-moyennes-frechet} doit tout poser sur place.
%     C'est l'ajout le plus rentable des cinq.
%  2. \cref{sec:context-riemann-orthogonal}: retraction polaire sur Stiefel et
%     notion de releve (lifting) comme logarithme approche. Actuellement seule
%     la retraction QR est mentionnee, en passant, dans
%     \cref{rem:context-riemann-stiefelvscone}.
%  3. \cref{sec:context-riemann-moyennes}: remarque sur la moyenne de Frechet
%     sur Stiefel et sur son inaccessibilite. C'est le pivot de
%     \cref{sec:spdnet-federe-obstacle}.
%  4. \cref{sec:context-riemann-notion}: proposition donnant les derivees
%     directionnelles de \expm{} et \logm{} (formule de Daleckii-Krein). Elles
%     servent dans toutes les propositions de \cref{sec:spdnet-batchnorm}, et
%     elles expliquent l'instabilite en $1/(\eig_\ivar-\eig_\ivarbis)$ de
%     \cref{prop:spdnet-gradevd}. Elles sont posees ici en
%     \cref{prop:spdnet-diffm} faute de mieux.
%  5. Figure: etendre l'experience qanat context_riemann_moyennes (aujourd'hui
%     arithmetique, log-euclidienne, Frechet) aux moyennes harmonique et GAH.
%     Elle servirait a la fois l'ajout 1 et
%     \cref{sec:spdnet-batchnorm-moyennes}, pour le cout d'un seul run.
% ====================================================================================


\section{Cadre commun: un réseau sous contrainte géométrique}
\label{sec:spdnet-cadre}

% ============================ TODO texte (~2 paragraphes) ===========================
% Ouvrir la section. Ce qu'il faut dire, dans cet ordre:
%  1. cette section n'est dans aucun des trois articles: elle est ecrite pour ce
%     memoire, et c'est elle qui fait du chapitre autre chose qu'une
%     juxtaposition de travaux --- meme role que
%     \cref{sec:detection-cadre} au chapitre precedent;
%  2. l'idee directrice: un reseau convolutif classique est une chaine
%     d'applications sur des espaces vectoriels; \SPDnet{} est la meme chaine
%     ou chaque espace est remplace par un cone
%     (\cref{sec:context-riemann-spd}) et chaque parametre par un point d'une
%     variete de Stiefel (\cref{def:context-riemann-stiefel}). Tout le reste du
%     chapitre consiste a payer le prix de cette substitution.
% ====================================================================================

\subsection{Les couches de \SPDnet}
\label{sec:spdnet-couches}

\begin{definition}[{Couche \BiMap}]
\label{def:spdnet-bimap}
  La couche \BiMap{} est la contrepartie de la convolution: à la transformation
  linéaire $\x\mapsto\Wlayer^\T\x$ d'un vecteur elle substitue la congruence
  \begin{equation}
    \MCov_\ilayer
    =
    \Wlayer_\ilayer^\T\,\MCov_{\ilayer-1}\,\Wlayer_\ilayer,
    \qquad
    \Wlayer_\ilayer\in\Stiefel{\dlayer{\ilayer-1}}{\dlayer{\ilayer}},
    \label{eq:spdnet-bimap}
  \end{equation}
  avec $\dlayer{\ilayer}\leq\dlayer{\ilayer-1}$ \sidecite[-0.5cm]{huang2017riemannian}.
  L'orthonormalité des colonnes de $\Wlayer_\ilayer$ garantit
  $\MCov_\ilayer\in\SPD{\dlayer{\ilayer}}$ dès que $\MCov_{\ilayer-1}\in\SPD{\dlayer{\ilayer-1}}$.
\end{definition}

\begin{definition}[{Couche \ReEig}]
\label{def:spdnet-reeig}
  La couche \ReEig{} est la contrepartie de la fonction d'activation \textsc{relu}:
  celle-ci s'appliquant terme à terme détruirait la structure SPD, elle est
  appliquée aux valeurs propres. Pour $\MCov_{\ilayer-1}=\Unit\eigset\Unit^\T$,
  \begin{equation}
    \MCov_\ilayer
    =
    \Unit\,\max\!\left(\reeigeps\Id{\dlayer{\ilayer}},\eigset\right)\Unit^\T,
    \label{eq:spdnet-reeig}
  \end{equation}
  où le seuil $\reeigeps>0$ n'est pas appris et est fixé avant l'entraînement.
\end{definition}

\begin{definition}[{Couche \LogEig}]
\label{def:spdnet-logeig}
  La couche \LogEig{} envoie la représentation dans un espace euclidien en la
  projetant sur l'espace tangent en l'identité
  (\cref{eq:context-riemann-explogspd}), puis en la vectorisant:
  \begin{equation}
    \MCov_{\Nlayer}
    =
    \logm\!\left(\MCov_{\Nlayer-1}\right),
    \qquad
    \hat{\classlabel}
    =
    \softmax\!\left(\MSoftmax\,\vech(\MCov_{\Nlayer})+\Vsoftmax\right),
    \label{eq:spdnet-logeig}
  \end{equation}
  avec $\MSoftmax\in\R^{\Nclass\times\dlayer{\Nlayer}(\dlayer{\Nlayer}+1)/2}$ et
  $\Vsoftmax\in\R^{\Nclass}$.
\end{definition}

\begin{definition}[{Architecture \SPDnet}]
\label{def:spdnet-architecture}
  Un réseau \SPDnet{} est la composition de $\Nlayer$ blocs \BiRe{}
  --- une couche \BiMap{} suivie d'une couche \ReEig{}, éventuellement d'une
  couche de normalisation par batch (\cref{def:spdnet-batchnorm}) ---, d'une
  couche \LogEig{} et d'une tête euclidienne. Ses paramètres sont
  \begin{equation}
    \netparam
    =
    \left(\Wlayer_1,\dots,\Wlayer_{\Nlayer},\MSoftmax,\Vsoftmax\right)
    \in
    \Stiefel{\dlayer{0}}{\dlayer{1}}\times\cdots\times
    \Stiefel{\dlayer{\Nlayer-1}}{\dlayer{\Nlayer}}\times\R^{\Nclass\times\frac{\dlayer{\Nlayer}(\dlayer{\Nlayer}+1)}{2}}\times\R^{\Nclass},
    \label{eq:spdnet-parametres}
  \end{equation}
  et l'application qu'il réalise est notée $\netfun_{\netparam}$.
\end{definition}

\begin{figure}[htb]
  \begin{center}
    \makebox[\textwidth][c]{% Architecture SPDnet, dessinee de haut en bas: c'est le seul sens qui laisse
% aux boites une taille lisible dans les 355pt du texte.
% Utilise dans la \cref{sec:spdnet-couches}.
% Styles et briques communs aux figures du \cref{ch:spdnet}.
% Les teintes reprennent celles de gfx/tikZ/riemann/style.tex pour que les
% figures des deux chapitres se lisent comme une meme famille.
% Le texte du memoire ne fait que 355pt (12,5 cm): toutes les figures de ce
% chapitre s'ecoulent donc de haut en bas, jamais de gauche a droite.
% Ce fichier peut etre \input plusieurs fois.
\definecolor{slayer}{RGB}{90,90,90}      % boite de couche, gris neutre
\definecolor{sspd}{RGB}{150,40,40}       % ce qui vit sur le cone SPD
\definecolor{sstiefel}{RGB}{40,70,140}   % ce qui vit sur Stiefel
\definecolor{seucl}{RGB}{50,110,120}     % ce qui vit dans un espace euclidien
\definecolor{sback}{RGB}{175,95,45}      % retropropagation
\tikzset{
  spdnet/.style={
    >={Stealth[length=5pt,width=4pt]},
    line join=round,
    font=\small,
    node distance=5mm,
  },
  % Boite de couche generique. La largeur est fixee pour que la colonne soit
  % alignee et que le texte respire.
  slayerbox/.style={
    draw=slayer!70, thick, rounded corners=2pt, align=center,
    inner sep=4pt, minimum height=9mm, minimum width=42mm, fill=white,
  },
  sstiefelbox/.style={slayerbox, draw=sstiefel!75, fill=sstiefel!7},
  splainbox/.style={slayerbox, fill=slayer!5},
  seuclbox/.style={slayerbox, draw=seucl!75, fill=seucl!7},
  smeanbox/.style={slayerbox, draw=sspd!75, fill=sspd!9},
  % Boite etroite, pour les colonnes doubles.
  snarrow/.style={slayerbox, minimum width=32mm},
  sarrow/.style={->, thick, slayer!85},
  sbackarrow/.style={->, thick, sback, dashed},
  sblock/.style={draw=slayer!45, thick, dashed, rounded corners=3pt, inner sep=5pt},
  slabel/.style={font=\footnotesize, inner sep=2pt},
  stiny/.style={font=\scriptsize, inner sep=1.5pt},
  % Accolade de regroupement a droite ou a gauche d'une colonne.
  sbrace/.style={decorate, decoration={brace, amplitude=4pt}, slayer!70},
}

\begin{tikzpicture}[spdnet]

  \node[splainbox] (in) {entrée $\MCov_0\in\SPD{\dlayer{0}}$};

  \node[sstiefelbox, below=of in] (bimap1) {\BiMap{} \ $\Wlayer_1\in\Stiefel{\dlayer{0}}{\dlayer{1}}$};
  \node[splainbox, below=of bimap1] (reeig1) {\ReEig{} \ seuil $\reeigeps$};
  \node[smeanbox, below=of reeig1] (bn1) {\BatchNorm{} \ $\Fmean,\ \Covbias_1$};

  \node[below=6mm of bn1] (dots) {$\vdots$};

  \node[sstiefelbox, below=6mm of dots] (bimapL) {\BiMap{} \ $\Wlayer_\Nlayer\in\Stiefel{\dlayer{\Nlayer-1}}{\dlayer{\Nlayer}}$};
  \node[splainbox, below=of bimapL] (reeigL) {\ReEig{} \ seuil $\reeigeps$};

  \node[splainbox, below=8mm of reeigL] (logeig) {\LogEig};
  \node[seuclbox, below=of logeig] (fc) {$\softmax$ \ $\MSoftmax,\Vsoftmax$};
  \node[splainbox, below=of fc] (out) {sortie $\hat{\classlabel}\in\R^{\Nclass}$};

  \foreach \a/\b in {in/bimap1, bimap1/reeig1, reeig1/bn1, bn1/dots, dots/bimapL,
                     bimapL/reeigL, reeigL/logeig, logeig/fc, fc/out}
    \draw[sarrow] (\a) -- (\b);
  \coordinate (mid) at ($(logeig.south)!0.5!(fc.north)$);
  \node[stiny, slayer, anchor=west] at ([xshift=2mm]mid) {$\vech$};

  % Regroupement des deux blocs BiRe.
  \node[sblock, fit=(bimap1)(bn1)] (b1) {};
  \node[stiny, slayer, rotate=90, anchor=south] at ($(b1.west)+(-1mm,0)$) {bloc \BiRe{} 1};
  \node[sblock, fit=(bimapL)(reeigL)] (bL) {};
  \node[stiny, slayer, rotate=90, anchor=south] at ($(bL.west)+(-1mm,0)$) {bloc \BiRe{} $\Nlayer$};

  % Colonne de droite: ou vivent les objets.
  \coordinate (x) at ([xshift=6mm]in.east);
  \draw[sspd!75, thick] ([yshift=1mm]in.north-|x) -- ([yshift=-1mm]logeig.north-|x);
  \node[stiny, sspd, anchor=west, align=left, text width=32mm]
    at ($([yshift=1mm]in.north-|x)!0.5!([yshift=-1mm]logeig.north-|x)+(2mm,0)$)
    {les représentations restent dans le cône $\SPD{\dlayer{\ilayer}}$};
  \draw[seucl!75, thick] ([yshift=-1mm]logeig.north-|x) -- ([yshift=-1mm]out.south-|x);
  \node[stiny, seucl, anchor=west, align=left, text width=32mm]
    at ($([yshift=-1mm]logeig.north-|x)!0.5!([yshift=-1mm]out.south-|x)+(2mm,0)$)
    {espace tangent en $\Id{\dlayer{\Nlayer}}$, puis $\R^{\Nclass}$};

\end{tikzpicture}
}
  \end{center}
  \caption{Architecture \SPDnet{} \sidecite[0cm]{huang2017riemannian}: une chaîne
    de $\Nlayer$ blocs \BiRe, suivie du passage au plan tangent par \LogEig{} et
    d'une tête euclidienne. % TODO texte: completer la legende.
  }
  \label{fig:spdnet-architecture}
\end{figure}

% ======================= TODO texte (remarque, ~10 lignes) ==========================
% Poser le label rem:spdnet-reeig-retrecissement. Deux lectures de \ReEig{} qui
% ne sont dans aucun des articles et qui rattachent le chapitre a la partie 2:
%  1. seuiller les valeurs propres par le bas est une *regularisation du
%     spectre*, de la meme famille que le retrecissement de
%     \cref{sec:context-covariance-regularized-estimation}: on paye un biais
%     pour garantir l'inversibilite et le conditionnement. La difference est
%     qu'ici la matrice a regulariser est produite par le reseau lui-meme;
%  2. donner les ordres de grandeur mesures dans \fullcite{PREP26b}: nombre de
%     conditionnement moyen de $9{,}1\times10^{5}$ (HDM05), $6{,}0\times10^{6}$
%     (HyperLeaf), $1{,}2\times10^{7}$ (Rices90). C'est ce qui rend le choix de
%     $\reeigeps$ critique et ce qui explique, en
%     \cref{sec:spdnet-batchnorm-resultats}, l'echec de la differentiation
%     automatique: les termes en $1/(\eig_\ivar-\eig_\ivarbis)$ de
%     \cref{prop:spdnet-gradevd} explosent exactement dans ce regime.
% Enchainer sur: \LogEig{} n'est pas une operation anodine non plus --- c'est le
% choix d'une geometrie, celle du plan tangent en $\Id{}$, donc une
% linearisation dont la validite depend de la dispersion des donnees.
% ====================================================================================

\subsection{Apprendre sous contrainte de Stiefel}
\label{sec:spdnet-apprentissage}

% ============================ TODO texte (~1 paragraphe) ============================
% Le strict minimum de vocabulaire, en une demi-page et pas davantage: on ne
% refait pas un cours d'apprentissage profond. Dire: perte d'entropie croisee,
% risque empirique, entrainement par batchs, et le fait que la seule
% specificite ici est que \cref{eq:spdnet-descente} est un pas riemannien au
% sens de \cref{eq:context-riemann-descente} pour les $\Wlayer_\ilayer$, et un
% pas euclidien ordinaire pour $\MSoftmax$ et $\Vsoftmax$. Preciser que, en
% pratique, les articles utilisent soit une retraction explicite, soit une
% parametrisation (trivialisation locale \sidecite{ablin2024infeasible},
% decomposition de Householder) qui ramene le probleme a une optimisation
% euclidienne --- et que ce choix est orthogonal a tout ce que dit le chapitre.
% ====================================================================================

Pour un jeu d'apprentissage $\{(\MCov_\iobs,\classlabel_\iobs)\}$ de
$\SPD{\dlayer{0}}\times\intrange{1}{\Nclass}$, le risque empirique associé à
l'entropie croisée s'écrit
\begin{equation}
  \loss(\netparam)
  =
  -\frac{1}{\Nobs}\sum_{\iobs=1}^{\Nobs}\sum_{\iclass=1}^{\Nclass}
  \mathbf{1}_{\{\classlabel_\iobs=\iclass\}}
  \log\left[\netfun_{\netparam}(\MCov_\iobs)\right]_\iclass,
  \label{eq:spdnet-risque}
\end{equation}
et le pas de descente s'écrit, pour les poids contraints,
\begin{equation}
  \Wlayer_\ilayer
  \gets
  \Retr_{\Wlayer_\ilayer}\!\left(-\lrate\,\Rgrad\loss(\Wlayer_\ilayer)\right),
  \qquad
  \Rgrad\loss(\Wlayer_\ilayer)
  =
  \projman_{\Wlayer_\ilayer}\!\left(\err{\loss}{\Wlayer_\ilayer}\right),
  \label{eq:spdnet-descente}
\end{equation}
avec $\projman$ la projection de \cref{eq:context-riemann-projstiefel}.

\subsection{Rétropropagation matricielle}
\label{sec:spdnet-backprop}

% ============================ TODO texte (~2 paragraphes) ===========================
% C'est le coeur technique de la section, et le denominateur commun aux trois
% travaux. A dire:
%  1. dans un reseau ordinaire, la retropropagation est une suite de produits
%     matrice-vecteur; ici les variables sont des matrices et les couches font
%     intervenir des decompositions spectrales, si bien que la regle de chaine
%     doit etre ecrite au niveau des differentielles, avec le produit scalaire
%     de Frobenius comme appariement \sidecite{ionescu2015matrix};
%  2. pourquoi ne pas se contenter de la differentiation automatique: elle
%     stocke le graphe complet, donc toutes les iterees d'un algorithme de point
%     fixe (\cref{eq:spdnet-geometrique-iteration}), d'ou une occupation memoire
%     qui croit avec le nombre d'iterations; et elle differentie une
%     decomposition spectrale terme a terme, avec les instabilites de
%     \cref{prop:spdnet-gradevd}. Les deux points sont mesures en
%     \cref{sec:spdnet-batchnorm-resultats}.
% ====================================================================================

\begin{proposition}[{Règle de chaîne matricielle}]
\label{prop:spdnet-chainrule}
  Soit une couche $\netfun_\ilayer$ transformant $\mat{X}_{\ilayer-1}$ en
  $\mat{X}_\ilayer=\netfun_\ilayer(\mat{X}_{\ilayer-1})$, et $\loss^{\ilayer}$ la
  perte vue depuis cette couche. L'erreur rétropropagée
  $\err{\loss^{\ilayer}}{\mat{X}_{\ilayer-1}}$ est l'unique matrice vérifiant,
  pour toute direction admissible $\diff\mat{X}_{\ilayer-1}$,
  \begin{equation}
    \trace{\err{\loss^{\ilayer}}{\mat{X}_{\ilayer-1}}^\T\diff\mat{X}_{\ilayer-1}}
    =
    \trace{\err{\loss^{\ilayer+1}}{\mat{X}_{\ilayer}}^\T\diff\mat{X}_{\ilayer}},
    \qquad
    \diff\mat{X}_\ilayer=\dirderiv\netfun_\ilayer(\mat{X}_{\ilayer-1})[\diff\mat{X}_{\ilayer-1}].
    \label{eq:spdnet-chainrule}
  \end{equation}
\end{proposition}

\begin{proposition}[{Dérivées directionnelles de $\expm$ et $\logm$}]
\label{prop:spdnet-diffm}
  Soit $\mat{X}=\Unit\eigset\Unit^\T$ la décomposition spectrale d'une matrice
  symétrique, $\eigset=\Diag(\eig_1,\dots,\eig_\Nvar)$, et $h$ une fonction
  scalaire appliquée au spectre. Pour toute direction symétrique $\tvec$,
  \begin{equation}
    \dirderiv h(\mat{X})[\tvec]
    =
    \Unit\left(\mat{G}\had\left(\Unit^\T\tvec\,\Unit\right)\right)\Unit^\T,
    \qquad
    \mat{G}_{\ivar\ivarbis}
    =
    \left\{
    \begin{array}{ll}
      \dfrac{h(\eig_\ivar)-h(\eig_\ivarbis)}{\eig_\ivar-\eig_\ivarbis}
        & \text{si } \eig_\ivar\neq\eig_\ivarbis,\\[8pt]
      h'(\eig_\ivar) & \text{sinon,}
    \end{array}
    \right.
    \label{eq:spdnet-daleckii}
  \end{equation}
  où $\had$ est le produit de Hadamard. Les cas $h=\exp$ et $h=\log$ donnent
  $\dirderiv\expm$ et $\dirderiv\logm$.
\end{proposition}

% ============================ TODO texte (~5 lignes) ================================
% Une phrase pour dire que \cref{prop:spdnet-diffm} est la formule de
% Daleckii-Krein, admise ici, et qu'elle aurait sa place au
% \cref{ch:context-riemann} (cf. le bloc TODO decision de l'introduction). Une
% phrase pour dire qu'elle est le seul ingredient manquant pour differentier
% toutes les couches spectrales du chapitre.
% ====================================================================================

\begin{proposition}[{Gradients des couches spectrales}]
\label{prop:spdnet-gradevd}
  Soit une couche s'écrivant $\mat{X}_\ilayer=\Unit\,h(\eigset)\,\Unit^\T$ à
  partir de $\mat{X}_{\ilayer-1}=\Unit\eigset\Unit^\T$. L'erreur rétropropagée est
  \begin{equation}
    \err{\loss^{\ilayer}}{\mat{X}_{\ilayer-1}}
    =
    2\,\Unit\left(\mat{P}^\T\had\symop\!\left(\Unit^\T\err{\loss^{\ilayer+1}\circ h}{\Unit}\right)\right)\Unit^\T
    +
    \Unit\,\Diag\!\left(\err{\loss^{\ilayer+1}\circ h}{\eigset}\right)\Unit^\T,
    \label{eq:spdnet-gradevd}
  \end{equation}
  où $\mat{P}_{\ivar\ivarbis}=(\eig_\ivar-\eig_\ivarbis)^{-1}$ pour
  $\ivar\neq\ivarbis$ et $\mat{P}_{\ivar\ivar}=0$
  \sidecite[-0.5cm]{ionescu2015matrix,brooks2019riemannian}. Pour la couche
  \ReEig{} de \cref{eq:spdnet-reeig},
  \begin{equation}
    \err{\loss^{\ilayer+1}\circ h}{\Unit}
    =
    2\,\symop\!\left(\err{\loss^{\ilayer+1}}{\mat{X}_{\ilayer}}\right)\Unit\max(\reeigeps\Id{},\eigset),
    \qquad
    \err{\loss^{\ilayer+1}\circ h}{\eigset}
    =
    \mat{Q}\,\Unit^\T\symop\!\left(\err{\loss^{\ilayer+1}}{\mat{X}_{\ilayer}}\right)\Unit,
    \label{eq:spdnet-gradreeig}
  \end{equation}
  avec $\mat{Q}=\Diag(\mathbf{1}_{\{\eig_1>\reeigeps\}},\dots)$; pour la couche
  \LogEig{} de \cref{eq:spdnet-logeig}, les mêmes expressions avec
  $\max(\reeigeps\Id{},\eigset)$ remplacé par $\log(\eigset)$ et $\mat{Q}$ par
  $\inv{\eigset}$.
\end{proposition}

% ======================= TODO texte (remarque, ~8 lignes) ===========================
% Poser le label rem:spdnet-sylvester. Deux points:
%  1. les equations de Sylvester du type $\mat{Z}\mat{S}+\mat{S}\mat{Z}=\mat{R}$
%     avec $\mat{S}\in\SPD{\Nvar}$ reviennent quatre fois dans ce chapitre
%     (\cref{prop:spdnet-batchnorm-backprop,prop:spdnet-grad-geo,prop:spdnet-grad-gah}).
%     Elles se resolvent en une decomposition spectrale de $\mat{S}$: dans la
%     base propre, $\mat{Z}_{\ivar\ivarbis}=\mat{R}_{\ivar\ivarbis}/(\std_\ivar+\std_\ivarbis)$,
%     ou les $\std$ sont les valeurs propres de $\mat{S}$. Le cout est celui
%     d'une decomposition spectrale, et la solution symetrique est unique
%     puisque $\std_\ivar+\std_\ivarbis>0$;
%  2. noter le contraste avec \cref{eq:spdnet-gradevd}: ici le denominateur est
%     une *somme* de valeurs propres, la ou la retropropagation spectrale fait
%     apparaitre une *difference*. C'est pourquoi les formules manuelles de ce
%     chapitre sont numeriquement plus sages que leur equivalent automatique.
% ====================================================================================

\subsection{Où se cachent les moyennes}
\label{sec:spdnet-feuille-route}

\begin{figure}[htb]
  \begin{center}
    \makebox[\textwidth][c]{% Feuille de route du chapitre: les trois endroits ou une moyenne euclidienne
% doit ceder la place a une moyenne de Frechet, la variete concernee, et la
% severite croissante de l'obstacle.
% Utilise dans la \cref{sec:spdnet-feuille-route}.
% Styles et briques communs aux figures du \cref{ch:spdnet}.
% Les teintes reprennent celles de gfx/tikZ/riemann/style.tex pour que les
% figures des deux chapitres se lisent comme une meme famille.
% Le texte du memoire ne fait que 355pt (12,5 cm): toutes les figures de ce
% chapitre s'ecoulent donc de haut en bas, jamais de gauche a droite.
% Ce fichier peut etre \input plusieurs fois.
\definecolor{slayer}{RGB}{90,90,90}      % boite de couche, gris neutre
\definecolor{sspd}{RGB}{150,40,40}       % ce qui vit sur le cone SPD
\definecolor{sstiefel}{RGB}{40,70,140}   % ce qui vit sur Stiefel
\definecolor{seucl}{RGB}{50,110,120}     % ce qui vit dans un espace euclidien
\definecolor{sback}{RGB}{175,95,45}      % retropropagation
\tikzset{
  spdnet/.style={
    >={Stealth[length=5pt,width=4pt]},
    line join=round,
    font=\small,
    node distance=5mm,
  },
  % Boite de couche generique. La largeur est fixee pour que la colonne soit
  % alignee et que le texte respire.
  slayerbox/.style={
    draw=slayer!70, thick, rounded corners=2pt, align=center,
    inner sep=4pt, minimum height=9mm, minimum width=42mm, fill=white,
  },
  sstiefelbox/.style={slayerbox, draw=sstiefel!75, fill=sstiefel!7},
  splainbox/.style={slayerbox, fill=slayer!5},
  seuclbox/.style={slayerbox, draw=seucl!75, fill=seucl!7},
  smeanbox/.style={slayerbox, draw=sspd!75, fill=sspd!9},
  % Boite etroite, pour les colonnes doubles.
  snarrow/.style={slayerbox, minimum width=32mm},
  sarrow/.style={->, thick, slayer!85},
  sbackarrow/.style={->, thick, sback, dashed},
  sblock/.style={draw=slayer!45, thick, dashed, rounded corners=3pt, inner sep=5pt},
  slabel/.style={font=\footnotesize, inner sep=2pt},
  stiny/.style={font=\scriptsize, inner sep=1.5pt},
  % Accolade de regroupement a droite ou a gauche d'une colonne.
  sbrace/.style={decorate, decoration={brace, amplitude=4pt}, slayer!70},
}

\begin{tikzpicture}[spdnet, node distance=7mm]

  \node[splainbox, minimum width=52mm] (amont)
    {\textbf{la donnée} \ --- \ mise en covariance};
  \node[smeanbox, minimum width=52mm, below=of amont] (dans)
    {\textbf{un batch} \ --- \ normalisation};
  \node[sstiefelbox, minimum width=52mm, below=of dans] (aval)
    {\textbf{les paramètres} \ --- \ agrégation};

  \draw[sarrow] (amont) -- (dans);
  \draw[sarrow] (dans) -- (aval);

  % Colonne de gauche: la moyenne concernee.
  \node[stiny, slayer, anchor=east, align=right, text width=30mm] at ([xshift=-4mm]amont.west)
    {$\CovPooled=\Tstack\centermat\Tstack^\T$\\sur $\SPD{\Nfilter}$};
  \node[stiny, slayer, anchor=east, align=right, text width=30mm] at ([xshift=-4mm]dans.west)
    {$\Fmean$ du batch\\sur $\SPD{\Nvar}$};
  \node[stiny, slayer, anchor=east, align=right, text width=30mm] at ([xshift=-4mm]aval.west)
    {moyenne des $\Wlayer_\ilayer$\\sur $\Stiefel{}{}$};

  % Colonne de droite: la section, et la severite de l'obstacle.
  \node[stiny, sspd, anchor=west, align=left, text width=27mm] at ([xshift=4mm]amont.east)
    {\cref{sec:spdnet-covpool}};
  \node[stiny, sspd, anchor=west, align=left, text width=27mm] at ([xshift=4mm]dans.east)
    {\cref{sec:spdnet-batchnorm}\\pas de forme close};
  \node[stiny, sspd, anchor=west, align=left, text width=27mm] at ([xshift=4mm]aval.east)
    {\cref{sec:spdnet-federe}\\pas même de logarithme};

  \node[stiny, slayer, below=5mm of aval, align=center, text width=95mm]
    {dans les trois cas, une moyenne de Fréchet (\cref{def:context-riemann-frechet}):
     il faut choisir une approximation, puis la différentier};

\end{tikzpicture}
}
  \end{center}
  \caption{Feuille de route du chapitre. % TODO texte: completer la legende.
  }
  \label{fig:spdnet-feuille-route}
\end{figure}

% ============================ TODO texte (~2 paragraphes) ===========================
% Section charniere, une page au plus. Ce qu'il faut poser:
%  1. les trois sections qui suivent traitent trois problemes qui n'ont, en
%     apparence, rien a voir: former la matrice d'entree, normaliser un batch,
%     agreger des poids entre clients. Enoncer qu'ils ont la meme forme: une
%     quantite qui, dans le cas euclidien, est une moyenne;
%  2. et que la reponse geometrique est chaque fois la meme --- la moyenne de
%     Frechet de \cref{def:context-riemann-frechet} --- avec le meme obstacle:
%     pas de forme close, donc un algorithme iteratif, donc un cout
%     inacceptable a chaque batch ou a chaque tour. Le degre de severite croit
%     d'ailleurs d'une section a l'autre: sur le cone, la moyenne de Frechet est
%     calculable mais chere (\cref{eq:context-riemann-frechetopt}); sur Stiefel,
%     le logarithme lui-meme n'a pas de forme close
%     (\cref{sec:spdnet-federe-obstacle});
%  3. annoncer les deux facons d'en sortir, qui structurent le reste du
%     chapitre: changer de geometrie pour une dont la moyenne est en forme close
%     (\cref{sec:spdnet-batchnorm-moyennes}), ou approcher les operateurs de la
%     bonne geometrie (\cref{sec:spdnet-federe-agregation}).
% ====================================================================================


\section{Former la matrice: la covariance comme couche apprise}
\label{sec:spdnet-covpool}

\begin{contributions}
  Cette section reprend \fullcite{JSTARS25}, ainsi que ses versions préliminaires
  \fullcite{IGARSS24} et \fullcite{GRETSI23b}.
\end{contributions}

% ============================ TODO texte (~2 paragraphes) ===========================
% Ouvrir la section. Elle doit rester courte --- environ le quart de
% \cref{sec:spdnet-batchnorm}. Ce qu'il faut dire:
%  1. le probleme applicatif en trois phrases: classifier des objets enfouis a
%     partir de vignettes d'hyperboles issues d'un radar geologique (\textsc{gpr}).
%     La difficulte n'est pas la forme de l'hyperbole en soi, c'est qu'elle
%     depend autant de l'objet que des conditions d'acquisition (frequence,
%     elevation du radar, nature et humidite du sol);
%  2. la these de la section: la representation du second ordre est un choix
%     *d'invariance*, et c'est ce qui la rend robuste au changement de
%     conditions. C'est la seule des trois sections du chapitre a mesurer cette
%     promesse, d'ou sa place en ouverture;
%  3. annoncer que la contribution methodologique tient en une proposition
%     (\cref{prop:spdnet-covpool-grad}) --- la covariance cesse d'etre une
%     caracteristique calculee en amont pour devenir une couche apprise --- et
%     que c'est cette prémisse dont \cref{sec:spdnet-batchnorm,sec:spdnet-federe}
%     ont besoin.
% ====================================================================================

\subsection{La couche de mise en covariance}
\label{sec:spdnet-covpool-couche}

\begin{definition}[{Couche \CovPool}]
\label{def:spdnet-covpool}
  Soit $\Tstack\in\R^{\Nfilter\times\Npix}$ la matrice regroupant les $\Nfilter$
  cartes de filtres d'un réseau convolutif, chacune aplatie en $\Npix$ pixels.
  La couche de mise en covariance produit
  \begin{equation}
    \CovPooled
    =
    \Tstack\,\centermat\,\Tstack^\T
    \in\SPD{\Nfilter},
    \qquad
    \centermat
    =
    \frac{1}{\Npix}\left(\Id{\Npix}-\frac{1}{\Npix}\vector{1}_{\Npix}\vector{1}_{\Npix}^\T\right),
    \label{eq:spdnet-covpool}
  \end{equation}
  c'est-à-dire la matrice de covariance empirique des $\Nfilter$ canaux, estimée
  sur les $\Npix$ positions spatiales.
\end{definition}

\begin{figure}[htb]
  \begin{center}
    \makebox[\textwidth][c]{% De la vignette GPR a la matrice SPD, de haut en bas. La retropropagation
% remonte a gauche: c'est elle qui fait de la covariance une couche apprise et
% non une caracteristique calculee en amont.
% Utilise dans la \cref{sec:spdnet-covpool-couche}.
% Styles et briques communs aux figures du \cref{ch:spdnet}.
% Les teintes reprennent celles de gfx/tikZ/riemann/style.tex pour que les
% figures des deux chapitres se lisent comme une meme famille.
% Le texte du memoire ne fait que 355pt (12,5 cm): toutes les figures de ce
% chapitre s'ecoulent donc de haut en bas, jamais de gauche a droite.
% Ce fichier peut etre \input plusieurs fois.
\definecolor{slayer}{RGB}{90,90,90}      % boite de couche, gris neutre
\definecolor{sspd}{RGB}{150,40,40}       % ce qui vit sur le cone SPD
\definecolor{sstiefel}{RGB}{40,70,140}   % ce qui vit sur Stiefel
\definecolor{seucl}{RGB}{50,110,120}     % ce qui vit dans un espace euclidien
\definecolor{sback}{RGB}{175,95,45}      % retropropagation
\tikzset{
  spdnet/.style={
    >={Stealth[length=5pt,width=4pt]},
    line join=round,
    font=\small,
    node distance=5mm,
  },
  % Boite de couche generique. La largeur est fixee pour que la colonne soit
  % alignee et que le texte respire.
  slayerbox/.style={
    draw=slayer!70, thick, rounded corners=2pt, align=center,
    inner sep=4pt, minimum height=9mm, minimum width=42mm, fill=white,
  },
  sstiefelbox/.style={slayerbox, draw=sstiefel!75, fill=sstiefel!7},
  splainbox/.style={slayerbox, fill=slayer!5},
  seuclbox/.style={slayerbox, draw=seucl!75, fill=seucl!7},
  smeanbox/.style={slayerbox, draw=sspd!75, fill=sspd!9},
  % Boite etroite, pour les colonnes doubles.
  snarrow/.style={slayerbox, minimum width=32mm},
  sarrow/.style={->, thick, slayer!85},
  sbackarrow/.style={->, thick, sback, dashed},
  sblock/.style={draw=slayer!45, thick, dashed, rounded corners=3pt, inner sep=5pt},
  slabel/.style={font=\footnotesize, inner sep=2pt},
  stiny/.style={font=\scriptsize, inner sep=1.5pt},
  % Accolade de regroupement a droite ou a gauche d'une colonne.
  sbrace/.style={decorate, decoration={brace, amplitude=4pt}, slayer!70},
}

\begin{tikzpicture}[spdnet]

  \node[splainbox] (img) {vignette $\imagegpr\in\R^{h\times w}$};
  \node[seuclbox, below=of img] (cnn) {$\ilayer$ premières couches d'un \layername{resnet}};
  \node[splainbox, below=of cnn] (stack) {cartes de filtres $\Tstack\in\R^{\Nfilter\times\Npix}$};
  \node[smeanbox, below=of stack] (cov) {\CovPool{} \ $\CovPooled=\Tstack\centermat\Tstack^\T$};
  \node[sstiefelbox, below=of cov] (spd) {\SPDnet{} \ $\Wlayer_1,\dots,\Wlayer_\Nlayer$};
  \node[seuclbox, below=of spd] (out) {\LogEig{} puis $\softmax$};

  \foreach \a/\b in {img/cnn, cnn/stack, stack/cov, cov/spd, spd/out}
    \draw[sarrow] (\a) -- (\b);

  \node[stiny, sspd, anchor=west, align=left, text width=34mm] at ([xshift=5mm]cov.east)
    {$\Nfilter$ variables pour $\Npix$ échantillons: le régime de la
     \cref{sec:context-covariance-regularized-estimation}};

  % Retropropagation: elle traverse la mise en covariance.
  \coordinate (l) at ([xshift=-6mm]img.west);
  \draw[sbackarrow] (out.west-|l) -- (cnn.west-|l);
  \node[stiny, sback, rotate=90, anchor=south, align=center] at ([xshift=-1.5mm]cov.west-|l)
    {la rétropropagation traverse\\la mise en covariance};

\end{tikzpicture}
}
  \end{center}
  \caption{% TODO texte: legende.
  }
  \label{fig:spdnet-covpool}
\end{figure}

% ======================= TODO texte (remarque, ~10 lignes) ==========================
% Poser le label rem:spdnet-covpool-regime. Point important et absent de
% l'article: \cref{eq:spdnet-covpool} est une matrice de covariance empirique
% estimee a partir de $\Npix$ echantillons pour $\Nfilter$ variables, c'est-a-dire
% exactement le regime dimensionnel discute en
% \cref{sec:context-covariance-regularized-estimation}. Donner les chiffres du
% papier: $\Nfilter=8\times32=256$ pour SRCNet, contre $\Npix=M_h\times M_w$
% de l'ordre de $38\times20=760$; le rapport $\Npix/\Nfilter$ est de l'ordre de 3,
% ce qui est faible. Conclure: c'est *pour cette raison* que la reduction de
% dimension par \BiMap{} et le seuil de \ReEig{} ne sont pas des raffinements
% mais des necessites --- sans eux, la matrice est trop mal conditionnee pour que
% \LogEig{} ait un sens. Rattacher a rem:spdnet-reeig-retrecissement.
% ====================================================================================

\subsection{Deux architectures: \layername{rcnet} et \layername{srcnet}}
\label{sec:spdnet-covpool-archi}

\begin{definition}[{\layername{rcnet} et \layername{srcnet}}]
\label{def:spdnet-rcnet}
  Les deux architectures diffèrent par la façon de former $\Tstack$ à partir des
  $\ilayer$ premières couches d'un réseau résiduel:
  \begin{itemize}
    \item \layername{rcnet}: seules les sorties de la dernière couche retenue
      sont utilisées, $\Tstack\in\R^{\Nfilter\times\Npix}$ avec
      $\Npix=h_\ilayer w_\ilayer$;
    \item \layername{srcnet}: les sorties des $\ilayer$ couches sont empilées
      après interpolation à une taille commune, ce qui donne
      $\Nfilter=\ilayer\cdot\bar{\Nfilter}$ cartes.
  \end{itemize}
  Dans les deux cas, $\CovPooled$ de \cref{eq:spdnet-covpool} alimente un réseau
  \SPDnet{} au sens de \cref{def:spdnet-architecture}.
\end{definition}

\begin{figure}[htb]
  \begin{center}
    \makebox[\textwidth][c]{% Les deux facons de former la matrice a partir d'un ResNet: RCNet ne garde que
% la derniere couche, SRCNet empile une partie des sorties de chacune des
% couches apres interpolation a une taille commune.
% Utilise dans la \cref{sec:spdnet-covpool-archi}.
% Styles et briques communs aux figures du \cref{ch:spdnet}.
% Les teintes reprennent celles de gfx/tikZ/riemann/style.tex pour que les
% figures des deux chapitres se lisent comme une meme famille.
% Le texte du memoire ne fait que 355pt (12,5 cm): toutes les figures de ce
% chapitre s'ecoulent donc de haut en bas, jamais de gauche a droite.
% Ce fichier peut etre \input plusieurs fois.
\definecolor{slayer}{RGB}{90,90,90}      % boite de couche, gris neutre
\definecolor{sspd}{RGB}{150,40,40}       % ce qui vit sur le cone SPD
\definecolor{sstiefel}{RGB}{40,70,140}   % ce qui vit sur Stiefel
\definecolor{seucl}{RGB}{50,110,120}     % ce qui vit dans un espace euclidien
\definecolor{sback}{RGB}{175,95,45}      % retropropagation
\tikzset{
  spdnet/.style={
    >={Stealth[length=5pt,width=4pt]},
    line join=round,
    font=\small,
    node distance=5mm,
  },
  % Boite de couche generique. La largeur est fixee pour que la colonne soit
  % alignee et que le texte respire.
  slayerbox/.style={
    draw=slayer!70, thick, rounded corners=2pt, align=center,
    inner sep=4pt, minimum height=9mm, minimum width=42mm, fill=white,
  },
  sstiefelbox/.style={slayerbox, draw=sstiefel!75, fill=sstiefel!7},
  splainbox/.style={slayerbox, fill=slayer!5},
  seuclbox/.style={slayerbox, draw=seucl!75, fill=seucl!7},
  smeanbox/.style={slayerbox, draw=sspd!75, fill=sspd!9},
  % Boite etroite, pour les colonnes doubles.
  snarrow/.style={slayerbox, minimum width=32mm},
  sarrow/.style={->, thick, slayer!85},
  sbackarrow/.style={->, thick, sback, dashed},
  sblock/.style={draw=slayer!45, thick, dashed, rounded corners=3pt, inner sep=5pt},
  slabel/.style={font=\footnotesize, inner sep=2pt},
  stiny/.style={font=\scriptsize, inner sep=1.5pt},
  % Accolade de regroupement a droite ou a gauche d'une colonne.
  sbrace/.style={decorate, decoration={brace, amplitude=4pt}, slayer!70},
}

\begin{tikzpicture}[spdnet, node distance=4mm]

  % ---------------------------- colonne gauche : RCNet ----------------------
  \begin{scope}
    \node[snarrow, seuclbox] (r1) {couche 1};
    \node[snarrow, seuclbox, below=of r1] (r2) {couche 2};
    \node[below=2mm of r2] (rd) {$\vdots$};
    \node[snarrow, seuclbox, below=2mm of rd] (rl) {couche $\ilayer$};
    \node[snarrow, splainbox, below=6mm of rl] (rt)
      {$\Tstack\in\R^{\Nfilter\times h_\ilayer w_\ilayer}$};
    \node[snarrow, smeanbox, below=of rt] (rc) {$\CovPooled\in\SPD{\Nfilter}$};

    \foreach \a/\b in {r1/r2, r2/rd, rd/rl} \draw[sarrow] (\a) -- (\b);
    \draw[sarrow] (rl) -- node[stiny, right] {$\Nfilter$ filtres} (rt);
    \draw[sarrow] (rt) -- (rc);

    \node[slabel, sstiefel, above=3mm of r1] {\textbf{\layername{rcnet}}};
    \node[stiny, slayer, below=3mm of rc, align=center, text width=34mm]
      {seule la dernière couche est utilisée};
  \end{scope}

  % ---------------------------- colonne droite : SRCNet ---------------------
  \begin{scope}[xshift=46mm]
    \node[snarrow, seuclbox] (s1) {couche 1};
    \node[snarrow, seuclbox, below=of s1] (s2) {couche 2};
    \node[below=2mm of s2] (sd) {$\vdots$};
    \node[snarrow, seuclbox, below=2mm of sd] (sl) {couche $\ilayer$};
    \node[snarrow, splainbox, below=6mm of sl] (st)
      {$\Tstack\in\R^{\ilayer\bar{\Nfilter}\times M_h M_w}$};
    \node[snarrow, smeanbox, below=of st] (sc) {$\CovPooled\in\SPD{\ilayer\bar{\Nfilter}}$};

    \foreach \a/\b in {s1/s2, s2/sd, sd/sl} \draw[sarrow] (\a) -- (\b);
    \draw[sarrow] (st) -- (sc);

    % Les prises sur chaque couche, ramenees vers l'empilement.
    \coordinate (tap) at ([xshift=8mm]s1.east);
    \foreach \n in {s1,s2,sl}
      \draw[sarrow, sback] (\n.east) -- ++(4mm,0) |- ([yshift=2mm]st.east);
    \node[stiny, sback, anchor=west, align=left, text width=24mm]
      at ([xshift=6mm,yshift=6mm]sl.east)
      {$\bar{\Nfilter}$ sorties par couche, interpolées};

    \node[slabel, sstiefel, above=3mm of s1] {\textbf{\layername{srcnet}}};
    \node[stiny, slayer, below=3mm of sc, align=center, text width=34mm]
      {les $\ilayer$ couches sont empilées};
  \end{scope}

\end{tikzpicture}
}
  \end{center}
  \caption{Les deux façons de former $\Tstack$ à partir des $\ilayer$ premières
    couches d'un réseau résiduel. % TODO texte: completer la legende --- dire ce
    % que chaque variante privilegie (profondeur des caracteristiques contre
    % richesse multi-echelle) et rappeler la contrainte memoire qui motive de ne
    % retenir que $\bar{\Nfilter}$ sorties par couche dans \layername{srcnet}.
  }
  \label{fig:spdnet-rcnet-srcnet}
\end{figure}

% ============================ TODO decision =========================================
% Voir \cref{fig:spdnet-rcnet-srcnet}. Les valeurs de reglage ($\ilayer=8$, $\bar{\Nfilter}=32$, $h=112$, $w=60$,
% interpolation a $38\times20$) sont dans l'article. Proposition: ne pas les
% reprendre ici, elles appartiennent au pdf joint, sauf celles qui servent
% l'argument du regime dimensionnel de rem:spdnet-covpool-regime. A confirmer.
% ====================================================================================

\subsection{Rétropropager jusqu'aux convolutions}
\label{sec:spdnet-covpool-backprop}

\begin{proposition}[{Gradient de la couche \CovPool}]
\label{prop:spdnet-covpool-grad}
  La couche \CovPool{} de \cref{eq:spdnet-covpool} n'a pas de paramètre à
  apprendre. L'erreur rétropropagée vers les cartes de filtres est
  \begin{equation}
    \err{\loss^{\ilayer}}{\Tstack}
    =
    2\,\centermat\,\Tstack^\T\,\err{\loss^{\ilayer+1}}{\CovPooled}.
    \label{eq:spdnet-covpool-grad}
  \end{equation}
\end{proposition}

% ============================ TODO texte (~1 paragraphe) ============================
% Ce que \cref{prop:spdnet-covpool-grad} change, et c'est le seul point
% methodologique de la section: la chaine de gradients est complete depuis la
% couche softmax jusqu'aux convolutions du reseau residuel. La matrice de
% covariance n'est donc plus une caracteristique calculee en amont, comme dans
% les approches par covariance pooling a poids figes \sidecite{li2017second,akodad2020ensemble},
% mais une couche dont l'amont est appris pour la tache. Enoncer que c'est la
% prémisse dont vivent les deux sections suivantes: sans elle, normaliser ou
% agreger a l'interieur du reseau n'aurait pas d'objet.
% ====================================================================================

\subsection{Classification d'objets enfouis en \textsc{gpr}}
\label{sec:spdnet-covpool-resultats}

% ==================== TODO texte + experiences (~2 pages) ===========================
% Trois experiences, toutes reprises de \fullcite{JSTARS25}, sur la base
% Geolithe (4 types d'objets enfouis, plusieurs configurations de frequence et
% d'elevation, sols secs et humides, sable et gravier). Convention du memoire:
% figures reprises d'un article -> citation, pas de macro \prov, les donnees
% n'etant pas distribuables. Protocole commun a rappeler une fois: 100 tirages
% de partition apprentissage/test, quantiles 5 / 50 / 95.
%
% FIGURE 1/3 --- precision de test en fonction de la fraction de donnees
%   d'apprentissage, pour CNN1, ResNet (pre-entraine et depuis zero), RCNet et
%   SRCNet. C'est la figure qui porte l'argument « peu de donnees ».
%
% FIGURE 2/3 --- precision de test en fonction du taux d'etiquettes erronees
%   dans l'ensemble d'apprentissage. Argument « robustesse au bruit
%   d'annotation ».
%
% FIGURE 3/3 --- les quatre scenarios de decalage entre apprentissage et test
%   (A: elevation du radar; B: frequence; C et D: nature du sol, gravier et
%   sable). C'est la figure la plus importante de la section: c'est elle qui
%   mesure ce que l'invariance geometrique promet. Decision a prendre: les
%   quatre scenarios en une figure a quatre panneaux, ou seulement A et C, les
%   deux autres renvoyant au pdf joint.
%
% Texte a ecrire autour: ne pas commenter scenario par scenario, mais degager la
% regularite --- l'ecart entre les modeles du second ordre et les autres *croit*
% avec la severite du decalage, alors qu'il est faible en conditions
% homogenes. C'est l'enonce a retenir pour le \cref{sec:spdnet-bilan}.
% ====================================================================================


\section{Normaliser un batch: quelle moyenne sur le cône?}
\label{sec:spdnet-batchnorm}

\begin{contributions}
  Cette section reprend \fullcite{PREP26b}.
\end{contributions}

% ============================ TODO texte (~2 paragraphes) ===========================
% Ouvrir la section. Ordre suggere:
%  1. rappeler ce que fait la normalisation par batch dans le cas euclidien
%     \sidecite{ioffe2015batch}: centrer et reduire chaque batch, puis
%     apprendre un biais qui rend a la couche la liberte de choisir la
%     distribution de sortie. C'est une des raisons pour lesquelles les reseaux
%     profonds s'entrainent;
%  2. sa transposition au cone \sidecite{brooks2019riemannian,kobler2022spd}
%     remplace la moyenne empirique par la moyenne geometrique. Enoncer les deux
%     limites qui motivent le travail: (i) la moyenne geometrique est iterative
%     et chere, (ii) la retropropagation a travers la couche n'avait pas ete
%     derivee et reposait sur la differentiation automatique;
%  3. annoncer le renversement de perspective de la section, qui est le point de
%     synthese du chapitre: le choix de la moyenne n'est pas un reglage, c'est le
%     choix d'une geometrie donc d'un modele statistique
%     (\cref{prop:spdnet-moyennes-frechet}). Le fait experimental que la moyenne
%     geometrique ne soit pas la meilleure devient alors un enonce sur les
%     donnees, pas une curiosite.
% Preciser aussi ce qui a ete volontairement laisse de cote: la normalisation de
% la dispersion \sidecite{kobler2022spd,9746629}, plus couteuse et instable pour
% des matrices mal conditionnees.
% ====================================================================================

\subsection{La couche de normalisation}
\label{sec:spdnet-batchnorm-couche}

\begin{definition}[{Normalisation par batch sur $\SPD{\Nvar}$}]
\label{def:spdnet-batchnorm}
  Soit $\{\MCov_\iobs\}_{\iobs=1}^{\Nobs}$ un batch de $\SPD{\Nvar}$ et $\Fmean$
  une moyenne de ce batch. La couche recentre le batch en l'identité puis lui
  applique un biais $\Covbias\in\SPD{\Nvar}$ appris:
  \begin{equation}
    \bnnorm_\iobs
    =
    \Fmean^{-1/2}\,\MCov_\iobs\,\Fmean^{-1/2},
    \qquad
    \bnout_\iobs
    =
    \Covbias\,\bnnorm_\iobs\,\Covbias,
    \label{eq:spdnet-batchnorm}
  \end{equation}
  de sorte que la moyenne du batch de sortie est $\Covbias^2$. En inférence, la
  moyenne du batch est remplacée par une moyenne globale $\Fmeanglob$ accumulée
  pendant l'apprentissage.
\end{definition}

% ======================= TODO texte (remarque, ~8 lignes) ===========================
% Poser le label rem:spdnet-biais-carre. Deux remarques de mise en place:
%  1. dans \sidecite{brooks2019riemannian,kobler2022spd} le biais est tel que la
%     moyenne du batch devienne $\Covbias$ et non $\Covbias^2$; le choix retenu
%     ici evite d'avoir a manipuler $\Covbias^{1/2}$ et allege toutes les
%     derivations qui suivent, sans rien changer a l'expressivite de la couche;
%  2. \cref{eq:spdnet-batchnorm} n'est autre que la congruence sous laquelle la
%     metrique affine invariante est invariante
%     (\cref{eq:context-riemann-invcongruence}). Autrement dit, *pour la
%     geometrie du cone, cette couche est une isometrie*: elle deplace les
%     donnees sans deformer leurs distances relatives. C'est une lecture qui
%     n'est dans aucun des articles et qui merite d'etre faite ici.
% ====================================================================================

\begin{property}[{Mise à jour de la moyenne globale}]
\label{prop:spdnet-momentum}
  La moyenne globale est mise à jour à chaque batch le long de la géodésique de
  la géométrie associée à la moyenne retenue. Pour la moyenne géométrique,
  \begin{equation}
    \Fmeanglob
    \gets
    \Fmeanglob\,\geodmean_{\momentum}\,\Fmean
    =
    (\Fmeanglob)^{1/2}\left((\Fmeanglob)^{-1/2}\Fmean(\Fmeanglob)^{-1/2}\right)^{\momentum}(\Fmeanglob)^{1/2},
    \label{eq:spdnet-momentum-geo}
  \end{equation}
  où $\momentum$ est le momentum et $\geodmean$ la géodésique de
  \cref{eq:context-riemann-geodspd}; les formules correspondantes pour les
  autres moyennes sont données en
  \cref{eq:spdnet-momentum-le,eq:spdnet-momentum-ari,eq:spdnet-momentum-harm,eq:spdnet-momentum-gah}.
\end{property}

% ======================= TODO texte (remarque, ~6 lignes) ===========================
% Poser le label rem:spdnet-momentum-geodesique. Le point a faire ressortir, qui
% n'est pas explicite dans l'article: la moyenne glissante par momentum, qui est
% une astuce d'implementation dans le cas euclidien, est ici *un point de la
% geodesique* entre l'ancienne moyenne globale et la moyenne du batch courant.
% Le momentum devient donc un objet geometrique, et chaque geometrie du chapitre
% fournit sa propre regle de mise a jour --- ce qui est le premier indice que le
% choix de la moyenne engage bien plus qu'un cout de calcul.
% ====================================================================================

\begin{figure}[p]
  \begin{center}
    \makebox[\textwidth][c]{% Couche de normalisation par batch: passe avant a gauche, passe arriere a
% droite, toutes deux de haut en bas. Ce qui rend la derivation non triviale est
% que l'entree agit sur la sortie par deux chemins, directement et par la
% moyenne: ils se rejoignent en bas a droite.
% Utilise dans la \cref{sec:spdnet-batchnorm-couche}.
% Styles et briques communs aux figures du \cref{ch:spdnet}.
% Les teintes reprennent celles de gfx/tikZ/riemann/style.tex pour que les
% figures des deux chapitres se lisent comme une meme famille.
% Le texte du memoire ne fait que 355pt (12,5 cm): toutes les figures de ce
% chapitre s'ecoulent donc de haut en bas, jamais de gauche a droite.
% Ce fichier peut etre \input plusieurs fois.
\definecolor{slayer}{RGB}{90,90,90}      % boite de couche, gris neutre
\definecolor{sspd}{RGB}{150,40,40}       % ce qui vit sur le cone SPD
\definecolor{sstiefel}{RGB}{40,70,140}   % ce qui vit sur Stiefel
\definecolor{seucl}{RGB}{50,110,120}     % ce qui vit dans un espace euclidien
\definecolor{sback}{RGB}{175,95,45}      % retropropagation
\tikzset{
  spdnet/.style={
    >={Stealth[length=5pt,width=4pt]},
    line join=round,
    font=\small,
    node distance=5mm,
  },
  % Boite de couche generique. La largeur est fixee pour que la colonne soit
  % alignee et que le texte respire.
  slayerbox/.style={
    draw=slayer!70, thick, rounded corners=2pt, align=center,
    inner sep=4pt, minimum height=9mm, minimum width=42mm, fill=white,
  },
  sstiefelbox/.style={slayerbox, draw=sstiefel!75, fill=sstiefel!7},
  splainbox/.style={slayerbox, fill=slayer!5},
  seuclbox/.style={slayerbox, draw=seucl!75, fill=seucl!7},
  smeanbox/.style={slayerbox, draw=sspd!75, fill=sspd!9},
  % Boite etroite, pour les colonnes doubles.
  snarrow/.style={slayerbox, minimum width=32mm},
  sarrow/.style={->, thick, slayer!85},
  sbackarrow/.style={->, thick, sback, dashed},
  sblock/.style={draw=slayer!45, thick, dashed, rounded corners=3pt, inner sep=5pt},
  slabel/.style={font=\footnotesize, inner sep=2pt},
  stiny/.style={font=\scriptsize, inner sep=1.5pt},
  % Accolade de regroupement a droite ou a gauche d'une colonne.
  sbrace/.style={decorate, decoration={brace, amplitude=4pt}, slayer!70},
}

\begin{tikzpicture}[spdnet, node distance=6mm]

  % ------------------------------- passe avant ------------------------------
  \begin{scope}
    \node[snarrow, splainbox] (in) {$\{\MCov_\iobs\}_{\iobs=1}^{\Nobs}$};
    \node[snarrow, smeanbox, below=of in] (mean) {moyenne $\Fmean$};
    \node[snarrow, splainbox, below=of mean] (norm)
      {$\bnnorm_\iobs=\Fmean^{-1/2}\MCov_\iobs\Fmean^{-1/2}$};
    \node[snarrow, sstiefelbox, below=of norm] (bias)
      {$\bnout_\iobs=\Covbias\,\bnnorm_\iobs\,\Covbias$};
    \node[snarrow, splainbox, below=of bias] (out) {$\{\bnout_\iobs\}_{\iobs=1}^{\Nobs}$};

    \draw[sarrow] (in) -- (mean);
    \draw[sarrow] (mean) -- (norm);
    \draw[sarrow] (norm) -- (bias);
    \draw[sarrow] (bias) -- (out);
    % le chemin direct, qui court-circuite la moyenne
    \draw[sarrow, rounded corners=3pt] (in.west) -- ++(-5mm,0)
      |- node[stiny, slayer, pos=0.25, rotate=90, anchor=south] {chemin direct} (norm.west);

    \node[slabel, slayer, above=3mm of in] {\textbf{passe avant}};
  \end{scope}

  % ------------------------------- passe arriere ----------------------------
  \begin{scope}[xshift=62mm]
    \node[snarrow, splainbox] (bout) {$\left\{\err{\loss}{\bnout_\iobs}\right\}$};
    \node[snarrow, sstiefelbox, below=of bout] (bbias) {$\err{\loss_{\bnout}}{\Covbias}$};
    \node[snarrow, splainbox, below=of bbias] (bnorm) {$\err{\loss_{\bnnorm}}{\MCov_\iobs}$};
    \node[snarrow, smeanbox, below=of bnorm] (bmean)
      {$\err{\loss_{\bnnorm}}{\Fmean}$ \ (Sylvester)};
    \node[snarrow, splainbox, below=of bmean] (bin)
      {$\err{\loss}{\MCov_\iobs}=\err{\loss_{\bnnorm}}{\MCov_\iobs}+\err{\loss_{\Fmean}}{\MCov_\iobs}$};

    \draw[sbackarrow] (bout) -- (bbias);
    \draw[sbackarrow] (bbias) -- (bnorm);
    \draw[sbackarrow] (bnorm) -- (bmean);
    \draw[sbackarrow] (bmean) -- (bin);
    \draw[sbackarrow, rounded corners=3pt] (bnorm.east) -- ++(5mm,0)
      |- node[stiny, sback, pos=0.25, rotate=90, anchor=north] {chemin direct} (bin.east);

    \node[slabel, sback, above=3mm of bout] {\textbf{passe arrière}};
  \end{scope}

\end{tikzpicture}
}
  \end{center}
  \caption{% TODO texte: legende. Insister sur les deux chemins de retour.
  }
  \label{fig:spdnet-batchnorm}
\end{figure}

\begin{algorithm}[htb]
  \caption{Normalisation par batch pour \SPDnet}
  \label{algo:spdnet-batchnorm}
  \begin{algorithmic}[1]
    \Statex\hspace*{-\algorithmicindent}\textbf{Apprentissage}
    \Require $\Nbatch$ batchs $\{\MCov_{\iobs,\ibatch}\}_{\iobs=1}^{\Nobs}$ de $\SPD{\Nvar}$, momentum $\momentum$
    \State Initialiser la moyenne globale $\Fmeanglob=\Id{\Nvar}$
    \For{$\ibatch\in\intrange{1}{\Nbatch}$}
      \State Calculer la moyenne du batch $\Fmean_\ibatch$
      \State Normaliser et biaiser selon \cref{eq:spdnet-batchnorm}
      \State Mettre à jour $\Fmeanglob$ selon \cref{eq:spdnet-momentum-geo}
    \EndFor
    \Ensure batchs transformés; paramètre appris $\Covbias$; moyenne estimée $\Fmeanglob$
    \Statex\hspace*{-\algorithmicindent}\textbf{Inférence}
    \Require $\MCov\in\SPD{\Nvar}$
    \Ensure $\Covbias\,(\Fmeanglob)^{-1/2}\,\MCov\,(\Fmeanglob)^{-1/2}\,\Covbias$
  \end{algorithmic}
\end{algorithm}

\subsection{Rétropropagation de la couche}
\label{sec:spdnet-batchnorm-backprop}

\begin{proposition}[{Rétropropagation de la normalisation par batch}]
\label{prop:spdnet-batchnorm-backprop}
  Soit $\{\err{\loss}{\bnout_\iobs}\}_{\iobs=1}^{\Nobs}$ les erreurs en
  sortie de la couche de \cref{def:spdnet-batchnorm}. L'erreur servant à mettre
  à jour le biais est
  \begin{equation}
    \err{\loss_{\bnout}}{\Covbias}
    =
    \sum_{\iobs=1}^{\Nobs}
    \left(
      \err{\loss}{\bnout_\iobs}\Covbias\,\bnnorm_\iobs
      +
      \bnnorm_\iobs\,\Covbias\err{\loss}{\bnout_\iobs}
    \right),
    \label{eq:spdnet-backprop-biais}
  \end{equation}
  et les erreurs de l'entrée se décomposent en deux contributions, l'une par la
  normalisation et l'autre par le calcul de la moyenne,
  \begin{equation}
    \err{\loss}{\MCov_\iobs}
    =
    \err{\loss_{\bnnorm}}{\MCov_\iobs}
    +
    \err{\loss_{\Fmean}}{\MCov_\iobs},
    \qquad
    \err{\loss_{\bnnorm}}{\MCov_\iobs}
    =
    \Fmean^{-1/2}\Covbias\err{\loss}{\bnout_\iobs}\Covbias\Fmean^{-1/2}.
    \label{eq:spdnet-backprop-entree}
  \end{equation}
  La seconde contribution s'obtient de $\err{\loss_{\bnnorm}}{\Fmean}$ par
  \begin{equation}
    \sum_{\iobs=1}^{\Nobs}\trace{\err{\loss_{\Fmean}}{\MCov_\iobs}\diff\MCov_\iobs}
    =
    \trace{\err{\loss_{\bnnorm}}{\Fmean}\diff\Fmean},
    \label{eq:spdnet-backprop-moyenne}
  \end{equation}
  où $\err{\loss_{\bnnorm}}{\Fmean}$ est l'unique solution symétrique de
  l'équation de Sylvester
  \begin{equation}
    \err{\loss_{\bnnorm}}{\Fmean}\Fmean^{1/2}
    +
    \Fmean^{1/2}\err{\loss_{\bnnorm}}{\Fmean}
    =
    \sum_{\iobs=1}^{\Nobs}
    \left(
      \err{\loss_{\bnnorm}}{\MCov_\iobs}\MCov_\iobs\Fmean^{-1/2}
      +
      \Fmean^{-1/2}\MCov_\iobs\err{\loss_{\bnnorm}}{\MCov_\iobs}
    \right).
    \label{eq:spdnet-backprop-sylvester}
  \end{equation}
\end{proposition}

% ============================ TODO texte (~1 paragraphe) ============================
% Commenter \cref{prop:spdnet-batchnorm-backprop}: la structure de la couche est
% telle que l'entree influence la sortie par deux chemins (directement par la
% normalisation, indirectement par la moyenne), d'ou les deux termes de
% \cref{eq:spdnet-backprop-entree} et la figure \ref{fig:spdnet-batchnorm}.
% Insister sur le fait que tout est explicite *sauf* le passage de
% \cref{eq:spdnet-backprop-moyenne}, qui depend de la moyenne choisie: c'est
% exactement le point ou la geometrie entre dans le calcul, et c'est l'objet de
% \cref{sec:spdnet-batchnorm-gradmoyennes}.
% ====================================================================================

\subsection{Cinq moyennes, cinq géométries}
\label{sec:spdnet-batchnorm-moyennes}

\begin{proposition}[{Les moyennes considérées sont des moyennes de Fréchet}]
\label{prop:spdnet-moyennes-frechet}
  Chacune des moyennes utilisées dans la couche de \cref{def:spdnet-batchnorm}
  est la moyenne de Fréchet (\cref{def:context-riemann-frechet}) du batch
  $\{\MCov_\iobs\}_{\iobs=1}^{\Nobs}$ pour une géométrie ou une divergence
  identifiée en \cref{pt:context}:
  \begin{equation}
    \begin{array}{rcl}
      \Fmean_{\meanlabel{G}}
        &=& \displaystyle\argmin_{\Fmean\in\SPD{\Nvar}}\ \sum_{\iobs}\geodist^2_{\meanlabel{AI}}(\Fmean,\MCov_\iobs),
        \\[10pt]
      \Fmean_{\meanlabel{LE}}
        &=& \displaystyle\argmin_{\Fmean\in\SPD{\Nvar}}\ \sum_{\iobs}\geodist^2_{\meanlabel{LE}}(\Fmean,\MCov_\iobs)
        \;=\; \expm\left(\frac{1}{\Nobs}\sum_{\iobs}\logm\MCov_\iobs\right),
        \\[10pt]
      \meanari
        &=& \displaystyle\argmin_{\Fmean\in\SPD{\Nvar}}\ \sum_{\iobs}\KL(\Fmean\,\|\,\MCov_\iobs)
        \;=\; \frac{1}{\Nobs}\sum_{\iobs}\MCov_\iobs,
        \\[10pt]
      \meanharm
        &=& \displaystyle\argmin_{\Fmean\in\SPD{\Nvar}}\ \sum_{\iobs}\KL(\MCov_\iobs\,\|\,\Fmean)
        \;=\; \left(\frac{1}{\Nobs}\sum_{\iobs}\inv{\MCov_\iobs}\right)^{-1},
        \\[10pt]
      \Fmean_{\meanlabel{GAH}}
        &=& \displaystyle\argmin_{\Fmean\in\SPD{\Nvar}}\ \sum_{\iobs}\left[\KL(\Fmean\,\|\,\MCov_\iobs)+\KL(\MCov_\iobs\,\|\,\Fmean)\right]
        \;=\; \meanari\,\geodmean_{1/2}\,\meanharm,
    \end{array}
    \label{eq:spdnet-moyennes-frechet}
  \end{equation}
  où $\geodist_{\meanlabel{AI}}$ et $\geodist_{\meanlabel{LE}}$ sont les distances
  affine invariante \cref{eq:context-riemann-distspd} et log-euclidienne, et où
  $\KL$ est la divergence de Kullback-Leibler entre lois normales centrées
  \cref{eq:context-geometry-kldivergence}
  \sidecite[-1cm]{moakher2012divergence,congedo2015approximate}.
\end{proposition}

% ==================== TODO texte (~3 paragraphes) --- LE POINT CENTRAL ==============
% C'est le passage le plus important du chapitre; le reste peut etre lu comme sa
% mise en oeuvre. Ce qu'il faut ecrire:
%  1. lire \cref{eq:spdnet-moyennes-frechet} comme un *prolongement de la
%     \cref{tab:context-riemann-metriques}*: la liste des moyennes n'est pas un
%     catalogue d'heuristiques numeriques, c'est une famille indexee par la
%     geometrie que l'on met sur le cone. Chaque ligne est une hypothese;
%  2. la lecture statistique, qui est ce qui transforme le resultat
%     experimental en enonce: la moyenne arithmetique est l'estimateur du
%     maximum de vraisemblance sous modele de Wishart, la moyenne harmonique son
%     equivalent sous modele de Wishart inverse, la moyenne geometrique
%     correspond a la distance de Rao du modele gaussien
%     (\cref{sec:context-geometry-gaussian}), et GAH est le compromis
%     symetrique entre les deux premieres. Dire alors ce qui est vraiment
%     mesure en \cref{sec:spdnet-batchnorm-resultats}: si GAH bat la moyenne
%     geometrique sur les trois jeux de donnees, ce n'est pas une curiosite
%     numerique, c'est que l'hypothese implicite de la moyenne geometrique n'est
%     pas la bonne pour ces donnees;
%  3. la remarque de cout, qui prepare \cref{sec:spdnet-batchnorm-cout}: la
%     ligne la mieux fondee du tableau est aussi la seule sans forme close.
%     C'est le meme obstacle qu'en \cref{sec:spdnet-feuille-route}, et la
%     reponse retenue ici est de *changer de geometrie*, la ou
%     \cref{sec:spdnet-federe} choisira d'approcher les operateurs.
% Ajouter une phrase sur la relation entre les deux premieres lignes: la moyenne
% log-euclidienne est exactement la premiere iteree de
% \cref{eq:spdnet-geometrique-iteration} initialisee en $\Id{\Nvar}$.
% ====================================================================================

La moyenne géométrique est obtenue par le point fixe \sidecite[0cm]{jeuris2012survey}
\begin{equation}
  \Fmean_{\iiter+1}
  =
  \Fmean_{\iiter}^{1/2}
  \expm\!\left(\frac{1}{\Nobs}\sum_{\iobs=1}^{\Nobs}
    \logm\!\left(\Fmean_{\iiter}^{-1/2}\MCov_\iobs\Fmean_{\iiter}^{-1/2}\right)\right)
  \Fmean_{\iiter}^{1/2},
  \qquad \Fmean_0=\Id{\Nvar},
  \label{eq:spdnet-geometrique-iteration}
\end{equation}
arrêté après un nombre fixé $\Niterm$ d'itérations. Les mises à jour de la
moyenne globale correspondant aux quatre autres moyennes s'écrivent
\begin{align}
  \Fmeanglob &\gets \expm\!\left((1-\momentum)\logm\Fmeanglob+\momentum\logm\Fmean_{\meanlabel{LE}}\right),
  \label{eq:spdnet-momentum-le}\\
  \Fmeanglob &\gets (1-\momentum)\,\Fmeanglob+\momentum\,\meanari,
  \label{eq:spdnet-momentum-ari}\\
  \Fmeanglob &\gets \left((1-\momentum)\inv{\Fmeanglob}+\momentum\inv{\meanharm}\right)^{-1},
  \label{eq:spdnet-momentum-harm}\\
  \Fmeanglob &\gets \left((1-\momentum)\Fmeanglob_{\meanari}+\momentum\,\meanari\right)
    \geodmean_{1/2}
    \left((1-\momentum)\inv{\Fmeanglob_{\meanharm}}+\momentum\inv{\meanharm}\right)^{-1}.
  \label{eq:spdnet-momentum-gah}
\end{align}

% ============================ TODO texte (~5 lignes) ================================
% Une phrase pour dire que \cref{eq:spdnet-momentum-gah} impose de maintenir
% separement les deux moyennes globales arithmetique et harmonique, ce qui est
% le seul surcout de memoire de GAH par rapport aux autres.
% ====================================================================================

\subsection{Rétropropagation des moyennes}
\label{sec:spdnet-batchnorm-gradmoyennes}

\begin{proposition}[{Rétropropagation de la moyenne géométrique}]
\label{prop:spdnet-grad-geo}
  Soit $\Fmean=\Fmean_{\Niterm}$ la sortie du point fixe
  \cref{eq:spdnet-geometrique-iteration} et $\err{\loss}{\Fmean}$ son erreur. En
  notant $\mat{\Gamma}_{\iobs,\iiter}=\Fmean_\iiter^{-1/2}\MCov_\iobs\Fmean_\iiter^{-1/2}$
  et
  \begin{equation}
    \mat{E}_\iiter
    =
    \dirderiv\expm\!\left(\tfrac{1}{\Nobs}\textstyle\sum_{\iobs}\logm\mat{\Gamma}_{\iobs,\iiter}\right)
    \left[\Fmean_\iiter^{1/2}\err{\loss^{(\Niterm-\iiter-1)}}{\Fmean_{\iiter+1}}\Fmean_\iiter^{1/2}\right],
    \label{eq:spdnet-grad-geo-E}
  \end{equation}
  les erreurs de l'entrée s'obtiennent en accumulant les contributions de toutes
  les itérations,
  \begin{equation}
    \err{\loss}{\MCov_\iobs}
    =
    \sum_{\iiter=0}^{\Niterm-1}
    \frac{1}{\Nobs}\,
    \Fmean_\iiter^{-1/2}\,
    \dirderiv\logm(\mat{\Gamma}_{\iobs,\iiter})\!\left[\mat{E}_\iiter\right]\,
    \Fmean_\iiter^{-1/2},
    \label{eq:spdnet-grad-geo}
  \end{equation}
  l'erreur sur l'itérée courante étant obtenue récursivement par
  $\err{\loss^{(\Niterm-\iiter)}}{\Fmean_\iiter}=\mat{M}_\iiter-\inv{\Fmean_\iiter}\mat{N}_\iiter\inv{\Fmean_\iiter}$,
  où $\mat{M}_\iiter$ et $\mat{N}_\iiter$ sont les uniques solutions symétriques des
  équations de Sylvester
  \begin{equation}
    \mat{M}_\iiter\Fmean_\iiter^{1/2}+\Fmean_\iiter^{1/2}\mat{M}_\iiter
    =
    \Fmean_\iiter^{-1/2}\Fmean_{\iiter+1}\err{\loss^{(\Niterm-\iiter-1)}}{\Fmean_{\iiter+1}}
    +
    \err{\loss^{(\Niterm-\iiter-1)}}{\Fmean_{\iiter+1}}\Fmean_{\iiter+1}\Fmean_\iiter^{-1/2},
    \label{eq:spdnet-grad-geo-M}
  \end{equation}
  \begin{equation}
    \mat{N}_\iiter\Fmean_\iiter^{-1/2}+\Fmean_\iiter^{-1/2}\mat{N}_\iiter
    =
    \frac{1}{\Nobs}\sum_{\iobs=1}^{\Nobs}
    \left(
      \MCov_\iobs\Fmean_\iiter^{-1/2}\dirderiv\logm(\mat{\Gamma}_{\iobs,\iiter})[\mat{E}_\iiter]
      +
      \dirderiv\logm(\mat{\Gamma}_{\iobs,\iiter})[\mat{E}_\iiter]\Fmean_\iiter^{-1/2}\MCov_\iobs
    \right).
    \label{eq:spdnet-grad-geo-N}
  \end{equation}
\end{proposition}

\begin{proposition}[{Rétropropagation de la moyenne log-euclidienne}]
\label{prop:spdnet-grad-le}
  Pour la moyenne log-euclidienne de \cref{eq:spdnet-moyennes-frechet},
  \begin{equation}
    \err{\loss}{\MCov_\iobs}
    =
    \frac{1}{\Nobs}\,\dirderiv\logm(\MCov_\iobs)\!\left[\mat{E}\right],
    \qquad
    \mat{E}
    =
    \dirderiv\expm\!\left(\tfrac{1}{\Nobs}\textstyle\sum_{\iobs}\logm\MCov_\iobs\right)\!\left[\err{\loss}{\Fmean}\right].
    \label{eq:spdnet-grad-le}
  \end{equation}
  C'est le cas $\Niterm=1$ de \cref{prop:spdnet-grad-geo}.
\end{proposition}

\begin{proposition}[{Rétropropagation des moyennes arithmétique et harmonique}]
\label{prop:spdnet-grad-arith-harm}
  Pour les moyennes $\meanari$ et $\meanharm$ de
  \cref{eq:spdnet-moyennes-frechet},
  \begin{equation}
    \err{\loss}{\MCov_\iobs}
    =
    \frac{1}{\Nobs}\err{\loss}{\meanari}
    \qquad\text{et}\qquad
    \err{\loss}{\MCov_\iobs}
    =
    \frac{1}{\Nobs}\,\inv{\MCov_\iobs}\,\meanharm\,\err{\loss}{\meanharm}\,\meanharm\,\inv{\MCov_\iobs}.
    \label{eq:spdnet-grad-arith-harm}
  \end{equation}
\end{proposition}

\begin{proposition}[{Rétropropagation de la moyenne GAH}]
\label{prop:spdnet-grad-gah}
  Pour $\Fmean_{\meanlabel{GAH}}=\meanari\geodmean_{1/2}\meanharm$,
  \begin{equation}
    \err{\loss}{\MCov_\iobs}
    =
    \frac{1}{\Nobs}\err{\loss}{\meanari}
    +
    \frac{1}{\Nobs}\inv{\MCov_\iobs}\meanharm\err{\loss}{\meanharm}\meanharm\inv{\MCov_\iobs},
    \label{eq:spdnet-grad-gah}
  \end{equation}
  où
  $\err{\loss}{\meanari}=\meanharm^{-1/2}\mat{M}\meanharm^{-1/2}$ et
  $\err{\loss}{\meanharm}=\meanari^{-1/2}\mat{N}\meanari^{-1/2}$, avec $\mat{M}$
  et $\mat{N}$ les uniques solutions symétriques des équations de Sylvester
  \begin{equation}
    \mat{M}\left(\meanharm^{-1/2}\meanari\meanharm^{-1/2}\right)^{1/2}
    +
    \left(\meanharm^{-1/2}\meanari\meanharm^{-1/2}\right)^{1/2}\mat{M}
    =
    \meanharm^{1/2}\err{\loss}{\Fmean}\meanharm^{1/2},
    \label{eq:spdnet-grad-gah-M}
  \end{equation}
  \begin{equation}
    \mat{N}\left(\meanari^{-1/2}\meanharm\meanari^{-1/2}\right)^{1/2}
    +
    \left(\meanari^{-1/2}\meanharm\meanari^{-1/2}\right)^{1/2}\mat{N}
    =
    \meanari^{1/2}\err{\loss}{\Fmean}\meanari^{1/2}.
    \label{eq:spdnet-grad-gah-N}
  \end{equation}
\end{proposition}

% ============================ TODO texte (~1 paragraphe) ============================
% Commenter la serie de propositions. Deux points seulement, il ne faut pas
% paraphraser les formules:
%  1. la hierarchie de complexite se lit directement: la moyenne arithmetique a
%     le gradient le plus simple qui soit, GAH demande deux equations de
%     Sylvester, la moyenne geometrique en demande deux *par iteration* du point
%     fixe --- et c'est ce dernier point qui explique le comportement memoire de
%     la differentiation automatique, qui doit stocker les $\Niterm$ iterees;
%  2. renvoyer aux materiels supplementaires de \fullcite{PREP26b} pour les
%     preuves, et indiquer que le code est disponible (depot
%     on_batchnormalization_for_spdnet).
% ====================================================================================

\subsection{Invariances et complexité}
\label{sec:spdnet-batchnorm-cout}

\begin{property}[{Invariances des moyennes}]
\label{prop:spdnet-invariances}
  Une moyenne est dite invariante par congruence si la moyenne de
  $\{\mat{A}\MCov_\iobs\mat{A}^\T\}$ vaut $\mat{A}\Fmean\mat{A}^\T$ pour toute
  $\mat{A}$ inversible, et invariante par inversion si la moyenne de
  $\{\inv{\MCov_\iobs}\}$ vaut $\inv{\Fmean}$. Les cinq moyennes de
  \cref{eq:spdnet-moyennes-frechet} se répartissent comme indiqué dans
  \cref{tab:spdnet-moyennes}: seules la moyenne géométrique et la moyenne GAH
  possèdent les deux.
\end{property}

\begin{table}[htb]
  \centering
  \caption{Les cinq moyennes de \cref{prop:spdnet-moyennes-frechet}: géométrie
    sous-jacente, lecture statistique, invariances et coût de calcul pour un
    batch de $\Nobs$ matrices de taille $\Nvar$.}
  \label{tab:spdnet-moyennes}
  \footnotesize
  \begin{tabular}{@{}llccc@{}}
    \toprule
    Moyenne & Géométrie / divergence & Congruence & Inversion & Coût \\
    \midrule
    Géométrique
      & affine invariante (Rao gaussien)
      & \checkmark & \checkmark
      & $\bigO(\Niterm\Nobs\Nvar^3)$ \\
    Log-euclidienne
      & log-euclidienne
      & $\times$ & \checkmark
      & $\bigO(\Nobs\Nvar^3)$ \\
    Arithmétique
      & $\KL$ à gauche (Wishart)
      & \checkmark & $\times$
      & $\bigO(\Nobs\Nvar^2)$ \\
    Harmonique
      & $\KL$ à droite (Wishart inverse)
      & \checkmark & $\times$
      & $\bigO(\Nobs\Nvar^3)$ \\
    \textsc{gah}
      & $\KL$ symétrisée
      & \checkmark & \checkmark
      & $\bigO(\Nobs\Nvar^3)$ \\
    \bottomrule
  \end{tabular}
\end{table}

% ============================ TODO texte (~1 paragraphe) ============================
% Lire \cref{tab:spdnet-moyennes} en trois phrases:
%  1. sur le plan des invariances, GAH est la seule alternative en forme close a
%     la moyenne geometrique qui ne perde rien --- c'est la prediction theorique
%     que les experiences vont confirmer;
%  2. la moyenne log-euclidienne perd l'invariance par congruence, ce qui est
%     coherent avec ce qui a ete dit en \cref{rem:context-riemann-LEbilan};
%  3. le facteur $\Niterm$ de la premiere ligne est le seul ecart de complexite
%     asymptotique, mais il se double du cout memoire de la differentiation
%     automatique, qui doit conserver les $\Niterm$ iterees. Preciser aussi que
%     la couche complete reste en $\bigO(\Nobs\Nvar^2+\Nvar^3)$ meme avec la
%     moyenne arithmetique, a cause de l'etape de normalisation elle-meme.
% ====================================================================================

\subsection{Coût mesuré et performances de classification}
\label{sec:spdnet-batchnorm-resultats}

% ==================== TODO texte + experiences (~3 pages) ===========================
% Deux volets, tous deux repris de \fullcite{PREP26b}. Le code etant public
% (depot on_batchnormalization_for_spdnet) et les jeux de donnees publics, ces
% resultats sont, contrairement a ceux de \cref{sec:spdnet-covpool-resultats},
% *reproductibles*. Decision a prendre: les rejouer via qanat pour porter les
% macros \prov, ou les reprendre de l'article avec citation. Le premier volet
% (simulations) serait peu couteux a rejouer; le second (trois jeux de donnees
% reels, recherche d'hyperparametres) ne l'est pas.
%
% VOLET 1 --- cout, sur donnees simulees. Matrices de nombre de conditionnement
%   $10^5$, comparaison differentiation manuelle / automatique pour la moyenne
%   geometrique et GAH.
%   FIGURE 1/3: temps et memoire GPU en fonction de la taille de matrice (8 a
%     512), batch de 64.
%   FIGURE 2/3: idem en fonction de la taille de batch (8 a 512), matrices 64.
%   FIGURE 3/3: idem en fonction de la profondeur (1 a 64 couches empilees).
%   Enonce a retenir: le gain memoire atteint un facteur 2 et *croit* avec les
%   trois parametres; le temps, lui, est comparable. C'est donc un argument de
%   memoire, pas de vitesse --- le dire franchement plutot que de survendre.
%
% VOLET 2 --- trois jeux de donnees reels, trois tableaux (F1, temps
%   d'entrainement, epoques a convergence, pic memoire GPU avec et sans
%   autograd), 5 graines, architecture a deux blocs BiRe.
%     - HDM05 \sidecite{muller2007mocap}: 2086 matrices $93\times93$, 117
%       classes, conditionnement moyen $9{,}1\times10^{5}$. GAH en tete, jusqu'a
%       5 % de F1 gagnes sur l'absence de normalisation; temps divise par deux au
%       moins par rapport a la moyenne geometrique;
%     - HyperLeaf \sidecite{Laprade_2024_CVPR}: 2410 images hyperspectrales de
%       feuilles de ble, 4 varietes, matrices $204\times204$, conditionnement
%       $6{,}0\times10^{6}$. Arithmetique, harmonique et GAH gagnent au moins
%       4 % de F1 sur la moyenne geometrique, 10 % sur l'absence de
%       normalisation; temps divise par 5, au prix de 10 a 20 epoques de plus;
%     - Rices90 \sidecite{8970539}: 10 des 90 especes, matrices $256\times256$,
%       conditionnement $1{,}2\times10^{7}$. Gain de 8 a 15 % de F1 par la
%       normalisation, peu d'ecart entre les moyennes, temps divise par 3,5.
%   Deux faits a mettre en avant et a relier explicitement:
%     (i) l'autograd *echoue* pour les moyennes geometrique et log-euclidienne
%         sur HyperLeaf et Rices90 --- ce n'est pas un gain marginal, c'est ce
%         qui rend ces methodes utilisables. Le rattacher au conditionnement des
%         donnees et a \cref{prop:spdnet-gradevd};
%     (ii) la moyenne geometrique, la mieux fondee theoriquement, n'est jamais la
%         meilleure. Renvoyer a \cref{prop:spdnet-moyennes-frechet} et conclure
%         comme un enonce de modelisation: sur ces donnees, l'hypothese
%         gaussienne sous-jacente a la geometrie affine invariante n'est pas la
%         bonne, et GAH, compromis symetrique entre Wishart et Wishart inverse,
%         l'est davantage. C'est la phrase que le chapitre doit produire.
% Signaler enfin l'artefact honnete de l'article: le temps plus long du modele
% sans normalisation vient de la taille de batch retenue par la recherche
% d'hyperparametres (8 contre davantage ailleurs), pas d'un effet de fond.
% ====================================================================================


\section{Agréger des poids: moyenner sur Stiefel}
\label{sec:spdnet-federe}

\begin{contributions}
  Cette section reprend \fullcite{EUSIPCO26b}.
\end{contributions}

% ============================ TODO texte (~2 paragraphes) ===========================
% Ouvrir la section. Ordre suggere:
%  1. le cadre federe en trois phrases \sidecite{mcmahan2017communication}:
%     apprendre un modele commun sans centraliser les donnees, en alternant
%     entrainements locaux et agregation sur un serveur. Motivation concrete du
%     travail: des enregistrements EEG repartis entre sites cliniques, qui ne
%     peuvent pas etre mis en commun;
%  2. l'obstacle, qui est le meme que celui de \cref{sec:spdnet-batchnorm} mais
%     un cran plus severe: la moyenne des poids est ici une moyenne sur la
%     variete de Stiefel, ou non seulement la moyenne de Frechet n'a pas de
%     forme close, mais le logarithme riemannien lui-meme n'en a pas
%     (\cref{prop:spdnet-federe-obstacle}). La reponse retenue ne sera donc pas
%     de changer de geometrie, mais d'approcher ses operateurs;
%  3. situer par rapport aux travaux existants sur l'apprentissage federe
%     riemannien \sidecite{li2022federated,huang2024riemannian,zhang2024nonconvex}:
%     ils traitent des varietes generales mais lient la methode a un optimiseur
%     local particulier (descente de gradient stochastique riemannienne). La
%     contribution ici est independante de l'optimiseur local, ce qui compte des
%     lors que les clients sont heterogenes.
% ====================================================================================

\subsection{Apprentissage fédéré et risque global}
\label{sec:spdnet-federe-probleme}

Soit $\clientset=\intrange{1}{\Nclient}$ l'ensemble des clients, le client
$\iclient$ disposant du jeu local
$\dataclient^{(\iclient)}=\{(\MCov_\iobs^{(\iclient)},\classlabel_\iobs^{(\iclient)})\}_{\iobs=1}^{|\dataclient^{(\iclient)}|}$
et du risque empirique local $\riskloc_\iclient(\netparam)$ défini comme en
\cref{eq:spdnet-risque}. L'objectif du serveur est
\begin{equation}
  \argmin_{\netparam}\ \frac{1}{\Nclient}\sum_{\iclient=1}^{\Nclient}\riskloc_\iclient(\netparam),
  \label{eq:spdnet-federe-risque}
\end{equation}
approché en $\Nround$ tours de communication. À chaque tour $\iround$, un
sous-ensemble $\selected_\iround\subset\clientset$ de $\Nselected$ clients est
tiré uniformément, chaque client sélectionné part de $\netparam_\iround$,
effectue $\Nepoch_\iclient$ époques locales et renvoie
$\netparam_\iround^{(\iclient)}$. Les paramètres euclidiens s'agrègent par
moyenne \sidecite[0cm]{mcmahan2017communication}
\begin{equation}
  \MSoftmax_{\iround+1}=\frac{1}{|\selected_\iround|}\sum_{\iclient\in\selected_\iround}\MSoftmax_\iround^{(\iclient)},
  \qquad
  \Vsoftmax_{\iround+1}=\frac{1}{|\selected_\iround|}\sum_{\iclient\in\selected_\iround}\Vsoftmax_\iround^{(\iclient)},
  \label{eq:spdnet-federe-euclidien}
\end{equation}
mais cette formule n'a pas de sens pour les poids $\Wlayer_{\ilayer,\iround}^{(\iclient)}$.

\subsection{L'obstacle: aucune moyenne de Fréchet accessible}
\label{sec:spdnet-federe-obstacle}

\begin{property}[{Inaccessibilité de la moyenne de Fréchet sur Stiefel}]
\label{prop:spdnet-federe-obstacle}
  La moyenne arithmétique de points de $\Stiefel{\Nvar}{k}$ n'appartient pas à
  $\Stiefel{\Nvar}{k}$ en général. La moyenne de Fréchet
  (\cref{def:context-riemann-frechet}) y est bien définie puisque la variété est
  compacte, mais son calcul par flot de Karcher requiert le logarithme
  riemannien et la distance géodésique, dont aucun n'admet de forme close sur
  $\Stiefel{\Nvar}{k}$: chacun demande lui-même une procédure itérative
  \sidecite[-0.5cm]{edelman1998geometry,bouchard2025beyond}. Son coût est donc
  prohibitif à chaque tour et pour chaque couche.
\end{property}

% ======================= TODO texte (remarque, ~8 lignes) ===========================
% Poser le label rem:spdnet-deux-obstacles. Mettre en regard explicite les deux
% obstacles du chapitre, c'est le moment ou le fil devient visible:
%  - sur le cone (\cref{sec:spdnet-batchnorm}), la moyenne de Frechet est
%    calculable, par un point fixe dont chaque iteration coute une
%    decomposition spectrale (\cref{eq:spdnet-geometrique-iteration}). Le
%    probleme est le prix, et la sortie a ete de changer de geometrie;
%  - sur Stiefel, le probleme n'est plus le prix mais la *disponibilite* des
%    operateurs eux-memes. Changer de geometrie ne sert a rien ici, il faut
%    approcher les operateurs: c'est ce que fournissent la retraction et le
%    releve. Renvoyer a \cref{rem:context-riemann-stiefelvscone}, qui avait deja
%    annonce que la compacite de Stiefel se paye par des exponentielles
%    couteuses.
% ====================================================================================

\subsection{Agréger par projection, ou par relèvement}
\label{sec:spdnet-federe-agregation}

\begin{definition}[{Agrégation par projection (\layername{projavg})}]
\label{def:spdnet-projavg}
  L'agrégation par projection remplace la moyenne de Fréchet par la projection
  de la moyenne arithmétique sur la variété:
  \begin{equation}
    \Wlayer_{\ilayer,\iround+1}
    =
    \polarf\!\left(\frac{1}{|\selected_\iround|}\sum_{\iclient\in\selected_\iround}\Wlayer_{\ilayer,\iround}^{(\iclient)}\right),
    \label{eq:spdnet-projavg}
  \end{equation}
  où $\polarf(\cdot)$ désigne le facteur orthogonal de la décomposition polaire,
  c'est-à-dire le point de $\Stiefel{\dlayer{\ilayer-1}}{\dlayer{\ilayer}}$ le
  plus proche au sens de la norme de Frobenius
  \sidecite[-0.5cm]{bouchard2025beyond}.
\end{definition}

\begin{definition}[{Relèvement et rétraction polaire}]
\label{def:spdnet-lift}
  La rétraction polaire et le relèvement associé sur
  $\Stiefel{\Nvar}{k}$ sont
  \begin{equation}
    \Retr_{\Wlayer}(\tvec)
    =
    \polarf(\Wlayer+\tvec),
    \qquad
    \Lift_{\Wlayer}(\bar{\Wlayer})
    =
    \projman_{\Wlayer}(\bar{\Wlayer}-\Wlayer),
    \label{eq:spdnet-lift}
  \end{equation}
  avec $\projman$ la projection sur l'espace tangent
  \cref{eq:context-riemann-projstiefel}. Elles jouent respectivement le rôle de
  l'exponentielle et du logarithme riemanniens
  \cref{eq:context-riemann-explogspd}, dont elles sont des approximations au
  premier ordre.
\end{definition}

\begin{definition}[{Agrégation par relèvement et rétraction (\layername{rlavg})}]
\label{def:spdnet-rlavg}
  Partant de la réécriture de la moyenne euclidienne
  $\bar{\Wlayer}=\exp_{\Wlayer}\!\left(\frac{1}{\Nselected}\sum_\iclient\log_{\Wlayer}(\Wlayer^{(\iclient)})\right)$,
  et en y substituant les approximations de \cref{def:spdnet-lift}, l'agrégation
  s'écrit
  \begin{equation}
    \Wlayer_{\ilayer,\iround+1}
    =
    \Retr_{\Wlayer_{\ilayer,\iround}}\!\left(
      \frac{1}{|\selected_\iround|}\sum_{\iclient\in\selected_\iround}
      \Lift_{\Wlayer_{\ilayer,\iround}}\!\left(\Wlayer_{\ilayer,\iround}^{(\iclient)}\right)
    \right).
    \label{eq:spdnet-rlavg}
  \end{equation}
\end{definition}

\begin{figure}[p]
  \begin{center}
    \makebox[\textwidth][c]{% ProjAvg: la moyenne arithmetique des poids locaux quitte la variete, on l'y
% ramene par le facteur orthogonal de la decomposition polaire.
% Figure dessinee pour ce memoire: elle n'existe dans aucun des articles.
% L'indice de couche est omis, il est le meme partout.
% Utilise dans la \cref{sec:spdnet-federe-agregation}.
% Styles et briques communs aux figures de geometrie riemannienne du
% \cref{ch:context-riemann}. Ce fichier peut etre \input plusieurs fois.
\definecolor{rmanifold}{RGB}{90,90,90}
\definecolor{rtangent}{RGB}{50,110,120}
\definecolor{rgeodesic}{RGB}{40,70,140}
\definecolor{rvector}{RGB}{175,95,45}
\tikzset{
  riemann/.style={
    >={Stealth[length=4.5pt,width=3.5pt]},
    line join=round,
    font=\footnotesize,
  },
  % La nappe qui represente la variete, centree sur l'origine.
  rsurface/.pic={
    \fill[rmanifold!7]
      (-3.2,0.15) .. controls (-2.95,1.75) and (-1.5,2.35) .. (0,2.35)
      .. controls (1.5,2.35) and (2.95,1.7) .. (3.2,0.15)
      .. controls (1.9,-0.7) and (-1.9,-0.7) .. (-3.2,0.15) -- cycle;
    \draw[rmanifold!70, thick]
      (-3.2,0.15) .. controls (-2.95,1.75) and (-1.5,2.35) .. (0,2.35)
      .. controls (1.5,2.35) and (2.95,1.7) .. (3.2,0.15);
    \draw[rmanifold!70, thick]
      (-3.2,0.15) .. controls (-1.9,-0.7) and (1.9,-0.7) .. (3.2,0.15);
  },
  % Le plan tangent, a placer en un point donne via un shift.
  rplane/.pic={
    \fill[rtangent!10, draw=rtangent!65]
      (-1.35,-0.42) -- (1.55,-0.88) -- (2.05,0.22) -- (-0.85,0.68) -- cycle;
  },
}

% Styles et briques communs aux figures du \cref{ch:spdnet}.
% Les teintes reprennent celles de gfx/tikZ/riemann/style.tex pour que les
% figures des deux chapitres se lisent comme une meme famille.
% Le texte du memoire ne fait que 355pt (12,5 cm): toutes les figures de ce
% chapitre s'ecoulent donc de haut en bas, jamais de gauche a droite.
% Ce fichier peut etre \input plusieurs fois.
\definecolor{slayer}{RGB}{90,90,90}      % boite de couche, gris neutre
\definecolor{sspd}{RGB}{150,40,40}       % ce qui vit sur le cone SPD
\definecolor{sstiefel}{RGB}{40,70,140}   % ce qui vit sur Stiefel
\definecolor{seucl}{RGB}{50,110,120}     % ce qui vit dans un espace euclidien
\definecolor{sback}{RGB}{175,95,45}      % retropropagation
\tikzset{
  spdnet/.style={
    >={Stealth[length=5pt,width=4pt]},
    line join=round,
    font=\small,
    node distance=5mm,
  },
  % Boite de couche generique. La largeur est fixee pour que la colonne soit
  % alignee et que le texte respire.
  slayerbox/.style={
    draw=slayer!70, thick, rounded corners=2pt, align=center,
    inner sep=4pt, minimum height=9mm, minimum width=42mm, fill=white,
  },
  sstiefelbox/.style={slayerbox, draw=sstiefel!75, fill=sstiefel!7},
  splainbox/.style={slayerbox, fill=slayer!5},
  seuclbox/.style={slayerbox, draw=seucl!75, fill=seucl!7},
  smeanbox/.style={slayerbox, draw=sspd!75, fill=sspd!9},
  % Boite etroite, pour les colonnes doubles.
  snarrow/.style={slayerbox, minimum width=32mm},
  sarrow/.style={->, thick, slayer!85},
  sbackarrow/.style={->, thick, sback, dashed},
  sblock/.style={draw=slayer!45, thick, dashed, rounded corners=3pt, inner sep=5pt},
  slabel/.style={font=\footnotesize, inner sep=2pt},
  stiny/.style={font=\scriptsize, inner sep=1.5pt},
  % Accolade de regroupement a droite ou a gauche d'une colonne.
  sbrace/.style={decorate, decoration={brace, amplitude=4pt}, slayer!70},
}

\begin{tikzpicture}[riemann, scale=1.62, every node/.style={transform shape}]

  \path (0,0) pic {rsurface};
  \node[rmanifold!85, font=\scriptsize] at (2.25,1.95)
    {$\Stiefel{\dlayer{\ilayer-1}}{\dlayer{\ilayer}}$};

  \coordinate (W1) at (-1.80,1.30);
  \coordinate (W2) at (-0.10,1.88);
  \coordinate (W3) at ( 1.55,1.22);
  \coordinate (Abar) at (-0.10,-1.30);
  \coordinate (Wnext) at (-0.10,0.95);

  \foreach \p in {W1,W2,W3} \draw[slayer!45, thin] (\p) -- (Abar);
  \fill[slayer!85] (Abar) circle (1.7pt);
  \node[slayer, font=\scriptsize, anchor=north, inner sep=3pt, align=center] at (Abar)
    {$\bar{\Wlayer}=\tfrac{1}{\Nselected}\sum_{\iclient}\Wlayer^{(\iclient)}$\\
     \textit{hors de la variété}};

  \foreach \p/\l in {W1/1, W2/2, W3/{\Nselected}}{
    \fill[sstiefel] (\p) circle (1.5pt);
    \node[sstiefel, font=\scriptsize, anchor=south, inner sep=2pt]
      at ([yshift=0.5pt]\p) {$\Wlayer^{(\l)}$};
  }

  \draw[sspd, very thick, ->] (Abar) -- (Wnext);
  \node[sspd, font=\scriptsize, anchor=west, inner sep=3.5pt] at ($(Abar)!0.5!(Wnext)$)
    {$\polarf$};

  \fill[sspd] (Wnext) circle (1.8pt);
  \node[sspd, font=\scriptsize, anchor=west, inner sep=3.5pt] at (Wnext)
    {$\Wlayer_{\iround+1}$};

\end{tikzpicture}
}
  \end{center}
  \caption[Agrégation par projection]{\layername{projavg}: la moyenne
    arithmétique des poids locaux quitte la variété, et l'on y revient par le
    facteur orthogonal de la décomposition polaire. % TODO texte: completer la
    % legende --- insister sur le fait que le point obtenu est le plus proche au
    % sens de Frobenius, et que l'iteree precedente n'intervient pas.
  }
  \label{fig:spdnet-projavg}
\end{figure}

\begin{figure}[p]
  \begin{center}
    \makebox[\textwidth][c]{% RLAvg: on releve les poids locaux dans l'espace tangent en l'iteree globale
% precedente, on y prend la moyenne --- qui a un sens, l'espace tangent etant
% vectoriel --- et l'on retourne sur la variete par retraction.
% Figure dessinee pour ce memoire: elle n'existe dans aucun des articles.
% Meme cadrage et meme echelle que agregation_projavg.tex, pour que les deux
% figures se comparent d'un coup d'oeil. L'indice de couche est omis.
% Utilise dans la \cref{sec:spdnet-federe-agregation}.
% Styles et briques communs aux figures de geometrie riemannienne du
% \cref{ch:context-riemann}. Ce fichier peut etre \input plusieurs fois.
\definecolor{rmanifold}{RGB}{90,90,90}
\definecolor{rtangent}{RGB}{50,110,120}
\definecolor{rgeodesic}{RGB}{40,70,140}
\definecolor{rvector}{RGB}{175,95,45}
\tikzset{
  riemann/.style={
    >={Stealth[length=4.5pt,width=3.5pt]},
    line join=round,
    font=\footnotesize,
  },
  % La nappe qui represente la variete, centree sur l'origine.
  rsurface/.pic={
    \fill[rmanifold!7]
      (-3.2,0.15) .. controls (-2.95,1.75) and (-1.5,2.35) .. (0,2.35)
      .. controls (1.5,2.35) and (2.95,1.7) .. (3.2,0.15)
      .. controls (1.9,-0.7) and (-1.9,-0.7) .. (-3.2,0.15) -- cycle;
    \draw[rmanifold!70, thick]
      (-3.2,0.15) .. controls (-2.95,1.75) and (-1.5,2.35) .. (0,2.35)
      .. controls (1.5,2.35) and (2.95,1.7) .. (3.2,0.15);
    \draw[rmanifold!70, thick]
      (-3.2,0.15) .. controls (-1.9,-0.7) and (1.9,-0.7) .. (3.2,0.15);
  },
  % Le plan tangent, a placer en un point donne via un shift.
  rplane/.pic={
    \fill[rtangent!10, draw=rtangent!65]
      (-1.35,-0.42) -- (1.55,-0.88) -- (2.05,0.22) -- (-0.85,0.68) -- cycle;
  },
}

% Styles et briques communs aux figures du \cref{ch:spdnet}.
% Les teintes reprennent celles de gfx/tikZ/riemann/style.tex pour que les
% figures des deux chapitres se lisent comme une meme famille.
% Le texte du memoire ne fait que 355pt (12,5 cm): toutes les figures de ce
% chapitre s'ecoulent donc de haut en bas, jamais de gauche a droite.
% Ce fichier peut etre \input plusieurs fois.
\definecolor{slayer}{RGB}{90,90,90}      % boite de couche, gris neutre
\definecolor{sspd}{RGB}{150,40,40}       % ce qui vit sur le cone SPD
\definecolor{sstiefel}{RGB}{40,70,140}   % ce qui vit sur Stiefel
\definecolor{seucl}{RGB}{50,110,120}     % ce qui vit dans un espace euclidien
\definecolor{sback}{RGB}{175,95,45}      % retropropagation
\tikzset{
  spdnet/.style={
    >={Stealth[length=5pt,width=4pt]},
    line join=round,
    font=\small,
    node distance=5mm,
  },
  % Boite de couche generique. La largeur est fixee pour que la colonne soit
  % alignee et que le texte respire.
  slayerbox/.style={
    draw=slayer!70, thick, rounded corners=2pt, align=center,
    inner sep=4pt, minimum height=9mm, minimum width=42mm, fill=white,
  },
  sstiefelbox/.style={slayerbox, draw=sstiefel!75, fill=sstiefel!7},
  splainbox/.style={slayerbox, fill=slayer!5},
  seuclbox/.style={slayerbox, draw=seucl!75, fill=seucl!7},
  smeanbox/.style={slayerbox, draw=sspd!75, fill=sspd!9},
  % Boite etroite, pour les colonnes doubles.
  snarrow/.style={slayerbox, minimum width=32mm},
  sarrow/.style={->, thick, slayer!85},
  sbackarrow/.style={->, thick, sback, dashed},
  sblock/.style={draw=slayer!45, thick, dashed, rounded corners=3pt, inner sep=5pt},
  slabel/.style={font=\footnotesize, inner sep=2pt},
  stiny/.style={font=\scriptsize, inner sep=1.5pt},
  % Accolade de regroupement a droite ou a gauche d'une colonne.
  sbrace/.style={decorate, decoration={brace, amplitude=4pt}, slayer!70},
}

\begin{tikzpicture}[riemann, scale=1.62, every node/.style={transform shape}]

  \path (0,0) pic {rsurface};
  \node[rmanifold!85, font=\scriptsize] at (2.25,1.95)
    {$\Stiefel{\dlayer{\ilayer-1}}{\dlayer{\ilayer}}$};

  \coordinate (W1) at (-1.80,1.30);
  \coordinate (W2) at (-0.10,1.88);
  \coordinate (W3) at ( 1.55,1.22);
  \coordinate (Wt) at (-0.30,0.15);
  \coordinate (Wnext) at ( 0.62,0.95);

  \path (Wt) pic[scale=0.98] {rplane};
  \node[rtangent!90, font=\scriptsize, anchor=east, inner sep=1pt]
    at ($(Wt)+(-1.42,0.44)$) {$\Tangent{\Wlayer_{\iround}}{\Stiefel{}{}}$};

  \coordinate (T1) at ($(Wt)+(-0.78,0.20)$);
  \coordinate (T2) at ($(Wt)+( 0.16,0.32)$);
  \coordinate (T3) at ($(Wt)+( 0.94,0.18)$);
  \coordinate (Tbar) at ($(Wt)+(0.11,0.24)$);
  \foreach \a/\b in {W1/T1, W2/T2, W3/T3}
    \draw[rtangent!75, thin, ->] (\a) -- (\b);
  \node[rtangent, font=\scriptsize, anchor=south east, inner sep=2pt] at ($(W1)!0.55!(T1)$)
    {$\Lift_{\Wlayer_\iround}$};

  \foreach \p/\l in {W1/1, W2/2, W3/{\Nselected}}{
    \fill[sstiefel] (\p) circle (1.5pt);
    \node[sstiefel, font=\scriptsize, anchor=south, inner sep=2pt]
      at ([yshift=0.5pt]\p) {$\Wlayer^{(\l)}$};
  }

  \fill (Wt) circle (1.5pt);
  \node[font=\scriptsize, anchor=north, inner sep=2.5pt] at (Wt) {$\Wlayer_{\iround}$};

  \fill[rtangent] (Tbar) circle (1.3pt);
  \node[rtangent, font=\scriptsize, anchor=north, inner sep=2.5pt, align=center]
    at ($(Tbar)+(0,-0.30)$)
    {$\tfrac{1}{\Nselected}\sum_{\iclient}\Lift_{\Wlayer_\iround}(\Wlayer^{(\iclient)})$};

  \draw[sspd, very thick, ->] (Tbar) .. controls ($(Tbar)+(0.42,0.28)$) .. (Wnext);
  \node[sspd, font=\scriptsize, anchor=west, inner sep=2.5pt] at ($(Tbar)+(0.30,0.42)$)
    {$\Retr_{\Wlayer_\iround}$};

  \fill[sspd] (Wnext) circle (1.8pt);
  \node[sspd, font=\scriptsize, anchor=west, inner sep=3pt] at (Wnext)
    {$\Wlayer_{\iround+1}$};

\end{tikzpicture}
}
  \end{center}
  \caption[Agrégation par relèvement et rétraction]{\layername{rlavg}: les poids
    locaux sont relevés dans l'espace tangent en l'itérée globale précédente, où
    la moyenne a un sens, puis ramenés sur la variété par rétraction.
    % TODO texte: completer la legende --- souligner le contraste avec la
    % \cref{fig:spdnet-projavg}: ici l'iteree precedente sert de point de
    % reference, ce qui oblige a la conserver.
  }
  \label{fig:spdnet-rlavg}
\end{figure}

\begin{algorithm}[htb]
  \caption{\layername{fedspdnet}}
  \label{algo:spdnet-fedspdnet}
  \begin{algorithmic}[1]
    \Require nombre de tours $\Nround$, clients par tour $\Nselected$, époques locales $\{\Nepoch_\iclient\}$
    \State Initialiser $\netparam_0=(\Wlayer_{1,0},\dots,\Wlayer_{\Nlayer,0},\MSoftmax_0,\Vsoftmax_0)$
    \For{$\iround=0,\dots,\Nround-1$}
      \State Tirer $\selected_\iround\subset\clientset$ avec $|\selected_\iround|=\Nselected$
      \For{\textbf{chaque} client $\iclient\in\selected_\iround$ \textbf{en parallèle}}
        \State Initialiser le modèle local avec $\netparam_\iround$
        \State Entraîner $\Nepoch_\iclient$ époques, obtenir $\netparam_\iround^{(\iclient)}$
      \EndFor
      \State Agréger $\MSoftmax_{\iround+1},\Vsoftmax_{\iround+1}$ par \cref{eq:spdnet-federe-euclidien}
      \For{$\ilayer=1,\dots,\Nlayer$}
        \State Agréger $\Wlayer_{\ilayer,\iround+1}$ par \cref{eq:spdnet-projavg} ou \cref{eq:spdnet-rlavg}
      \EndFor
    \EndFor
    \Ensure $\netparam_{\Nround}$
  \end{algorithmic}
\end{algorithm}

\subsection{Équivalence des deux schémas}
\label{sec:spdnet-federe-equivalence}

\begin{proposition}[{Équivalence des deux agrégations}]
\label{prop:spdnet-federe-equivalence}
  Les agrégations \cref{eq:spdnet-projavg} et \cref{eq:spdnet-rlavg} ont la même
  complexité par tour,
  $\bigO(\Nselected\,\dlayer{\ilayer-1}\dlayer{\ilayer}+\dlayer{\ilayer-1}\dlayer{\ilayer}^2)$.
  Lorsque les poids locaux restent proches de l'itérée globale, c'est-à-dire
  $\Wlayer^{(\iclient)}_{\ilayer,\iround}=\Wlayer_{\ilayer,\iround}+\bigO(\epsilon)$
  uniformément en $\iclient$, elles coïncident à l'ordre $\bigO(\epsilon^2)$ près.
\end{proposition}

% ============================ TODO texte (~1 paragraphe) ============================
% Justifier \cref{prop:spdnet-federe-equivalence} en deux lignes de calcul dans
% le texte (le releve est lineaire en $\bar{\Wlayer}-\Wlayer$, donc sa moyenne
% est le projete tangent de la moyenne arithmetique centree; la retraction
% polaire de ce vecteur est le projete polaire de la moyenne arithmetique, a
% l'ordre deux pres). Enoncer la consequence pratique: \layername{projavg} est
% recommandee, car elle a des constantes plus faibles et n'exige pas de
% conserver l'iteree globale precedente. Renvoyer aux courbes de
% \cref{sec:spdnet-federe-resultats} qui montrent des trajectoires superposees.
% ====================================================================================

\subsection{Classification \textsc{eeg} multi-sites}
\label{sec:spdnet-federe-resultats}

% ==================== TODO texte + experiences (~2 pages) ===========================
% Resultats repris de \fullcite{EUSIPCO26b}, sur deux jeux de donnees publics
% d'imagerie motrice EEG issus de MOABB \sidecite{chevallier2024largest}:
%   - Weibo2014: 10 sujets, 60 voies, 7 classes, ~80 essais par classe;
%   - PhysionetMI: 106 sujets (3 exclus), 64 voies, 4 classes, ~23 essais.
% Entree: matrice de covariance empirique de chaque essai centre, filtre en
% [8, 32] Hz. Un seul bloc BiRe. Reference euclidienne: EEGnet
% \sidecite{lawhern2018eegnet} federe par FedAvg. Federation simulee par
% partition des sujets, deux par client ($\Nclient=5$ et $53$), $\Nround=150$
% tours, $\Nepoch_\iclient=2$, participation totale et a moitie.
%
% FIGURE 1/2 --- F1 de validation en fonction du tour, pour les deux jeux, avec
%   \layername{projavg}, \layername{rlavg} et FedEEGnet. Deux enonces: les
%   courbes des deux agregations se superposent (c'est la verification
%   empirique de \cref{prop:spdnet-federe-equivalence}), et FedSPDnet atteint son
%   plateau en moins de tours.
%
% TABLEAU 1/1 --- F1 final, centralise / participation totale / participation a
%   moitie, avec le nombre de parametres. Chiffres a reprendre: Weibo2014,
%   SPDnet 51,7 en centralise et 43,3 en federe (perte 8,4) contre EEGnet 50,7 et
%   39,9 (perte 10,8); PhysionetMI, SPDnet 43,1 et 39,5 (perte 3,6) contre EEGnet
%   50,8 et 38,3 (perte 12,5). Nombres de parametres: 4 715 contre 5 303, et
%   2 452 contre 3 284.
%
% Texte a ecrire autour, et il faut etre honnete sur le point suivant: sur
% PhysionetMI, SPDnet est *battu* en centralise et le reste ne doit pas etre
% presente autrement. Ce qui est remarquable est ailleurs: c'est la quantite de
% performance *conservee* sous federation, et la stabilite a la participation
% partielle. Formuler l'hypothese explicative avec la prudence de l'article ---
% la contrainte de Stiefel restreint l'espace des parametres et limite la derive
% des clients --- en indiquant que ce n'est pas demontre.
% Mentionner enfin la limite annoncee dans l'article: garder les donnees chez le
% client ne garantit pas la confidentialite, les poids communiques restant
% attaquables; la confidentialite differentielle sur variete est ouverte, et
% c'est un des points repris au \cref{ch:perspectives}.
% ====================================================================================


\section{Bilan: trois substitutions, un même arbitrage}
\label{sec:spdnet-bilan}

\begin{table}[htb]
  \centering
  \caption{Les trois substitutions du chapitre. Dans chaque cas, l'objet
    géométriquement correct est une moyenne de Fréchet inaccessible, et la
    contribution consiste à choisir l'approximation puis à en dériver le
    gradient.}
  \label{tab:spdnet-bilan}
  \footnotesize
  \begin{tabular}{@{}p{2.2cm}p{2.6cm}p{2.6cm}p{2.6cm}p{2.6cm}@{}}
    \toprule
    & Opération standard & Contrepartie géométrique & Obstacle & Réponse retenue \\
    \midrule
    \cref{sec:spdnet-covpool}
      & caractéristique du premier ordre
      & covariance apprise sur $\SPD{\Nfilter}$
      & régime dimensionnel, conditionnement
      & réduction \BiMap{} et seuil \ReEig, gradient explicite \\
    \cref{sec:spdnet-batchnorm}
      & moyenne empirique du batch
      & moyenne de Fréchet sur $\SPD{\Nvar}$
      & pas de forme close pour la métrique affine invariante
      & changer de géométrie: moyennes en forme close \\
    \cref{sec:spdnet-federe}
      & moyenne des poids
      & moyenne de Fréchet sur $\Stiefel{}{}$
      & ni logarithme ni distance en forme close
      & approcher les opérateurs: rétraction et relèvement \\
    \bottomrule
  \end{tabular}
\end{table}

% ============================ TODO texte (~4 paragraphes) ===========================
% Bilan du chapitre. Ce qu'il doit produire, dans cet ordre:
%  1. relire \cref{tab:spdnet-bilan} comme une progression et non comme une
%     liste: le meme probleme se pose trois fois, de plus en plus severement, et
%     les deux sorties possibles --- changer de geometrie, ou approcher ses
%     operateurs --- sont exhibees l'une apres l'autre;
%  2. l'enseignement de modelisation, qui est ce que ce chapitre apporte de plus
%     specifique: le choix d'une moyenne est un choix de modele
%     (\cref{prop:spdnet-moyennes-frechet}), et il est mesurable. C'est aussi la
%     ou le memoire depasse les articles, qui presentent ce choix comme un
%     arbitrage cout/performance;
%  3. l'enseignement numerique: deriver a la main n'est pas une coquetterie. La
%     ou la differentiation automatique echoue (\cref{sec:spdnet-batchnorm-resultats}),
%     c'est la structure de la formule --- sommes de valeurs propres au
%     denominateur plutot que differences --- qui sauve le calcul;
%  4. l'ouverture vers le \cref{ch:perspectives}: la moyenne ponderee que calcule
%     un mecanisme d'attention est une moyenne de Karcher ponderee, donc la
%     quatrieme occurrence du meme probleme, et la question du cout s'y pose dans
%     les memes termes (\cref{sec:perspectives-transformers}). Mentionner aussi
%     la mise a disposition logicielle: anotherspdnet pour les couches de
%     \cref{sec:spdnet-covpool}, le depot de \fullcite{PREP26b} pour la
%     normalisation, et \fullcite{PREP26} (SPD Learn) pour la suite.
% ====================================================================================

% ============ TODO texte (~1 page) --- LES SIX MOUVEMENTS DE CHATELAIN =============
% A INSERER AVANT le point 4 ci-dessus, qui doit rester la derniere chose du
% chapitre. Meme cahier des charges qu'en \cref{sec:detection-synthese}: les
% points 1 a 3 couvrent la science, ce bloc couvre ce qu'un jury lit.
%
%  (a) NOMMER LES PERSONNES, AVEC LES DATES ET LES STATUTS. Ce chapitre est
%      celui ou la distinction compte le plus:
%        - Jafuno: these (\fullcite{JSTARS25});
%        - Pautrel: stage M2 en 2025 puis post-doc (\fullcite{EUSIPCO26b});
%        - Gallet: COLLEGUE MCF, collaboration entre pairs, EXPLICITEMENT PAS
%          un encadrement (\fullcite{PREP26b}). Le dire est a votre credit:
%          un jury verifie ce genre de distinction.
%
%  (b) L'APPLICATION QUI A UNIFIE LE TRAVAIL: le GPR ouvre le chapitre et reste
%      court parce que c'est la seule des trois sections qui MESURE la
%      robustesse que l'invariance geometrique promet. Le rappeler ici.
%
%  (c) CE QUI A ETE CONSTRUIT, EN CLAIR, sans equations.
%
%  (d) POINTER LES ARTICLES REPRODUITS AVEC LEUR FOLIO:
%      \paperpage{JSTARS25} (p. 191), \paperpage{PREP26b} (p. 205),
%      \paperpage{EUSIPCO26a} (p. 215), \paperpage{EUSIPCO26b} (p. 221).
%
%  (e) DIRE QUE LE TRAVAIL EXCEDE LES ARTICLES JOINTS, et citer le reste
%      (dont \fullcite{IGARSS24}, liste dans l'encadre sans pdf).
%
%  (f) DIRE OU CELA CONTINUE MAINTENANT, nommement.
%
% Puis 3 ou 4 PERSPECTIVES PROPRES AU CHAPITRE en \paragraph{Titre.} gras.
% ====================================================================================

% ====================================================================================
% Articles reproduits. \fullcite{EUSIPCO26a} (ARMAGNAC) est la version conference
% du travail de normalisation par lots dont \fullcite{PREP26b} est la version
% journal; son pdf est dans papers/ mais l'encadre d'ouverture ne le liste pas.
% \fullcite{IGARSS24}, lui, est liste dans l'encadre sans pdf disponible.
% ====================================================================================
\joinedpapers{%
  Les versions publiées des travaux dont rend compte ce chapitre sont reproduites
  ci-après, dans l'ordre où le chapitre les mobilise.

  \joinedpaper[JSTARS25]{\citeauthor{JSTARS25}, \citetitle{JSTARS25}}{papers/Classification_of_Buried_Objects_From_Ground_Penetrating_Radar_Images_by_Using_Second-Order_Deep_Learning_Models.pdf}
  \joinedpaper[PREP26b]{\citeauthor{PREP26b}, \citetitle{PREP26b}}{papers/1-s2.0-S0165168426003063-main.pdf}
  \joinedpaper[EUSIPCO26a]{\citeauthor{EUSIPCO26a}, \citetitle{EUSIPCO26a}}{papers/Eusipco26_BN_pub.pdf}
  \joinedpaper[EUSIPCO26b]{\citeauthor{EUSIPCO26b}, \citetitle{EUSIPCO26b}}{papers/2604.22494v1.pdf}
}
